\documentclass[10pt,a4paper]{article}

% ----------------- Paquetes -------------------%
\usepackage[T1]{fontenc}
\usepackage[utf8]{inputenc}
\usepackage[a4paper,margin=2.5cm]{geometry}
\usepackage{mathptmx}             
\usepackage{changepage} 
\usepackage{setspace}
\usepackage[spanish]{babel}
\usepackage[labelfont=bf,labelsep=space]{caption}
\usepackage{amssymb}
\usepackage{amsmath}
\usepackage{ebgaramond}
\usepackage{placeins}
\usepackage{etoolbox}
\usepackage{setspace}
\usepackage{parskip} 
\usepackage{tcolorbox}
\usepackage{needspace}
\usepackage{xparse}
\usepackage{ocgx2}
\usepackage{hyperref}
\usepackage{hypcap}
\usepackage{soul}\sethlcolor{yellow}% resaltado amarillo: \hl{texto} (Panel TCEE, ⌃H)


%%%%%%%%%%%%%%%%%%%%%%%%%%%%%%%%%%%%%%%%%%%%%%%%%%%%%%%%%%%%%%%%
% --- Fija primero tus valores globales "buenos" ---
\setstretch{1.2}                 % Interlineado global
\setlength{\parskip}{0.7\baselineskip}
\setlength{\parindent}{0pt}
\newlength{\GlobalParskip}
\setlength{\GlobalParskip}{\parskip}

\newcounter{eqnum}[subsection]
\renewcommand{\theeqnum}{\arabic{eqnum}}
\newcounter{alphaitem}
\newcounter{numitem}

%% ------------------- Recuadros----------------------%%
\newcommand{\puntosclave}[1]{%
  \begin{tcolorbox}[
    colframe=black,
    title={Puntos clave},
    before upper={\setlength{\parindent}{0.5cm}\setlength{\parskip}{0.5\baselineskip}}
  ]
  #1
  \end{tcolorbox}
}

\newcommand{\modificaciones}[1]{
  \begin{tcolorbox}[
    colback=white,
    colframe=red,
    title={Modificaciones a realizar},
    before upper={\setlength{\parindent}{0.5cm}\setlength{\parskip}{0.5\baselineskip}}
  ]
  #1
  \end{tcolorbox}
}

%% ------------------- CITA ----------------------%%
% \begin{cita}[Autor][Año][Obra] texto \end{cita}: cita sangrada y, debajo a la derecha, «AUTOR (Año), Obra». Los tres datos son opcionales.
% En el Panel TCEE: escribir \cita e Intro (el cursor va al texto; Tab pasa a Autor, Año y Obra).
\NewDocumentEnvironment{cita}{ O{} O{} O{} +b }{%
  \begin{adjustwidth}{2cm}{0pt}
  \raggedright
  #4\par
  \raggedleft
  \if\relax\detokenize{#1#2#3}\relax
    % sin firma
  \else
    #1%
    \if\relax\detokenize{#2}\relax\else\if\relax\detokenize{#1}\relax\else\space\fi(#2)\fi
    \if\relax\detokenize{#3}\relax\else\if\relax\detokenize{#1#2}\relax\else, \fi\textit{#3}\fi
  \fi
  \end{adjustwidth}
}{}



%% ------------------- Listas----------------------%%
    % --- Lista alfabética ---
    \newenvironment{la}
      {\begin{list}{\alph{alphaitem})}{%
          \usecounter{alphaitem}%
          \setlength{\leftmargin}{1cm}%
          \setlength{\parskip}{3pt}% 
          \setlength{\itemsep}{0pt}% ← espacio entre ítems
          \setlength{\parsep}{0pt}%  ← párrafos dentro de un ítem
      }}
      {\end{list}}
      
    % --- Lista numérica ---
    \newenvironment{lnum}
      {\begin{list}{\arabic{numitem}.}{%
          \usecounter{numitem}%
          \setlength{\leftmargin}{1cm}%
          \setlength{\parskip}{3pt}% 
          \setlength{\itemsep}{0pt}% ← espacio entre ítems
          \setlength{\parsep}{0pt}%  ← párrafos dentro de un ítem
      }}
      {\end{list}}
      
% ------------------- Imágenes----------------------%
    \graphicspath{{Img/}}% Rutas por defecto 
    % --- Imagen adaptativa ---
    % Registros de dimensión definidos una sola vez
    \newdimen\imgcap
    \newdimen\imgmin
    \newdimen\imgmax
    \newdimen\imgremain
    \newdimen\imgh
    \newdimen\imgneed
    
    \newcommand{\imagenfit}[3]{%
      \addvspace{10pt}% separación antes
      \par\begingroup
        % Parámetros de composición (asignaciones, no \newdimen)
        \imgcap=1\baselineskip            % alto estimado de título+fuente
        \imgmin=0.15\textheight
        \imgmax=0.15\textheight
        \imgremain=\pagegoal
        \ifdim\pagegoal=\maxdimen \imgremain=0pt\fi
        \advance\imgremain by -\pagetotal
        \advance\imgremain by -\imgcap
        % Decisión de altura destino
        \ifdim\imgremain<\imgmin
          \newpage
          \imgh=0.20\textheight
        \else
          \imgh=\imgremain
          \ifdim\imgh>\imgmax \imgh=\imgmax \fi
        \fi
        % Reservar espacio para evitar cortes del bloque completo
        \imgneed=\imgh \advance\imgneed by \imgcap
        \Needspace{\imgneed}
        % Cuerpo
        \noindent\begin{minipage}{\textwidth}
          \centering
          \captionof{figure}{\textemdash\ #2}\par\vspace{0.3em}% título arriba
          \includegraphics[width=\linewidth,height=\imgh,keepaspectratio]{\detokenize{#1}}\par
          \vspace{0.3em}\small\textit{Fuente: }#3% fuente abajo
        \end{minipage}%
      \endgroup
      \par\addvspace{10pt}%
    }

%------------ Bloque de RECUADRO ESPECIAL -------------%
    \newenvironment{specialblock}[1]{%
      \begin{quote} % crea sangría a izquierda y derecha
      \small % texto más pequeño
      \noindent\textbf{#1}\\[-0.3em] % título en negrita
      \rule{\linewidth}{0.4pt}\\[0.6em]}{
      \end{quote}}
      

%-------------------------- Author Font -----------------------%
    \newcommand{\authorfont}[1]{{\fontfamily{EBGaramond-TLF}\selectfont #1}}

%-------------------------- Anexos -----------------------%
    \makeatletter
    \newcounter{anexo}
    \renewcommand{\theanexo}{\Roman{anexo}}
    \newcommand{\anexo}[1]{%
      \clearpage
      \phantomsection
      \refstepcounter{anexo}%
      \noindent\textbf{Anexo \theanexo.- }#1\par\medskip
      \addcontentsline{stc}{section}{Anexo \theanexo.- #1}%
    }
    \makeatother

    %-------------------------- Lista de figuras (personal) -----------------------%
\makeatletter
\newcommand{\MyListOfFigures}{%
  \clearpage
  \phantomsection
  {\centering\bfseries\Large Índice de figuras\par}%
  \medskip
  % Entrada en nuestro índice de contenido (stc)
  \addcontentsline{stc}{section}{Índice de figuras}%
  % Recuperación del listado (sin cabecera propia)
  \begingroup
    \parindent=0pt\parskip=2pt
    \@starttoc{lof}%
  \endgroup
}
\makeatother

%-------------------------- Bibliografía -----------------------%
\makeatletter
% Contador para las referencias
\newcounter{myref}
% Cabecera de la bibliografía, entra en el índice 'stc'
\newcommand{\MyBibliography}{%
  \clearpage
  \phantomsection
  {\centering\bfseries\Large Bibliografía y referencias\par}%
  \medskip
  \addcontentsline{stc}{section}{Bibliografía y referencias}%
  \setcounter{myref}{0}%
}
\newcommand{\MyBibItem}[4]{%
  \refstepcounter{myref}%
  \noindent[\themyref]\ #1.\ \emph{#2}. #3. #4\par
}
\makeatother

%--------- Establecer cuatro niveles de índice --------%
    \setcounter{tocdepth}{4}   
    \setcounter{secnumdepth}{4}% numera hasta \paragraph

%%%%%%%%%%%%%%%%%%%%%%%%%%%%%%%%

\makeatletter

    %--------- Interlineado Global --------%
    \edef\GlobalBaseLineStretch{\baselinestretch}
    
    %--------- Espaciado de seccinones --------%
    \renewcommand\section{\@startsection{section}
      {1}{\z@}
      {2.5ex \@plus 1ex \@minus .2ex} % espacio antes
      {1.5ex \@plus .2ex}             % espacio después
      {\normalfont\Large\bfseries}    % formato del título
    }
    \renewcommand\subsection{\@startsection{subsection}{2}{\z@}
      {3ex \@plus 0.8ex \@minus .2ex}
      {1.5ex \@plus .2ex}
      {\normalfont\large\bfseries}
    }
    \renewcommand\subsubsection{\@startsection{subsubsection}{3}{\z@}
      {3ex \@plus 0.5ex \@minus .2ex}
      {1ex \@plus .2ex}
      {\normalfont\normalsize\bfseries}
    }
    \renewcommand\paragraph{\@startsection{paragraph}{4}{\z@}%
    {-2ex \@plus -0.5ex \@minus -.2ex}% espacio antes
    {0.8ex \@plus .2ex}% espacio después
    {\normalfont\normalsize\itshape}}% ← cursiva en lugar de negrita

    %--------- Ecuaciones --------%   
    % Variables globales para almacenar títulos y notas
    \gdef\leq@current@section{}%
    \gdef\leq@current@subsection{}%
    \gdef\leq@last@printed@section{}%
    \gdef\leq@last@printed@subsection{}%
    
    % Comando para añadir nota (invisible en el cuerpo)
    \newcommand{\eqnote}[1]{%
      \addcontentsline{leq}{leqnote}{#1}%
    }
    
    % Comando para ecuaciones
    \newcommand{\eqblock}[2]{%
      % Verificar si necesitamos imprimir el título de sección
      \ifx\leq@current@section\leq@last@printed@section\else
        \addcontentsline{leq}{leqsection}{\leq@current@section}%
        \global\let\leq@last@printed@section\leq@current@section
        \gdef\leq@last@printed@subsection{}% Reset subsección al cambiar sección
      \fi
      % Verificar si necesitamos imprimir el título de subsección
      \ifx\leq@current@subsection\leq@last@printed@subsection\else
        \ifx\leq@current@subsection\@empty\else
          \addcontentsline{leq}{leqsubsection}{\leq@current@subsection}%
          \global\let\leq@last@printed@subsection\leq@current@subsection
        \fi
      \fi
      % Ecuación
      \refstepcounter{eqnum}%
      \begin{equation}
        #1\tag{E.\theeqnum}
      \end{equation}%
      \addcontentsline{leq}{leqt}{[E.\theeqnum] -- #2}%
      \addcontentsline{leq}{leqm}{#1}%
    }
    
    %%% ----- Índice de ecuaciones ----------%%%
    % Tamaño para []
    \newcommand{\leqIndexFont}{\normalsize}
    \newcommand{\leqTitleFont}{\tiny}
    \newcommand{\leqNoteFont}{\tiny}
    % Cabecera de sección 
    \newcommand\l@leqsection[2]{%
      \addvspace{25pt}% Mayor separación antes de nueva sección
      \noindent\leqIndexFont
      \makebox[\linewidth]{\rule{.38\linewidth}{0.4pt}\ \ \textbf{  \MakeUppercase{#1}}\ \ \rule{.38\linewidth}{0.4pt}}%
      \par\vspace{-10pt}
    }
    % Cabecera de subsección 
    \newcommand\l@leqsubsection[2]{%
      \par\addvspace{15pt}%
      \noindent\leqIndexFont #1
      \par\vspace{-7pt}
    }
    % Título de ecuación 
    \newcommand\l@leqt[2]{%
    \par\addvspace{10pt}%
      \par\noindent\leqTitleFont #1\par
    }
    % Miniatura de ecuación
    \newcommand\l@leqm[2]{%
      \vspace{0.3\baselineskip}%
      \par\scriptsize\noindent\hspace{1.5em}\(#1\)\par%
    }
    % Nota de ecuación 
    \newcommand\l@leqnote[2]{%
      \vspace{0.2\baselineskip}%
      \par\noindent\leqNoteFont\hspace{1.5em}\textbf{Nota.--} #1\par\vspace{0.5\baselineskip}%
    }
    % Generación del índice
    \newcommand{\mylistofequations}{%
      \clearpage
      \phantomsection
      % Título visual de la lista (ajústalo a tu gusto):
      {\centering\bfseries\Large Índice de ecuaciones\par}
      % Entrada en nuestro índice 'stc':
      \addcontentsline{stc}{section}{Índice de ecuaciones}%
      % Llamada a tu lista real (usa el nombre exacto de tu macro):
      \@starttoc{leq}%
    }
    
    %--------- Captura de títulos de sección --------%
    \let\leq@original@section\section
    \let\leq@original@subsection\subsection
    % Flag para saber si estamos en una sección asterisco
    \newif\ifLeqStarSection
    \LeqStarSectionfalse
    
    % Redefinir \section CON CONTROL DE ASTERISCO
    \renewcommand{\section}{%
      \@ifstar{\leq@section@star}{\@dblarg\leq@section@nostar}%
    }
    \def\leq@section@star#1{%
      \global\LeqStarSectiontrue
      \gdef\leq@current@section{#1}%
      \gdef\leq@current@subsection{}%
      \leq@original@section{#1}%
      \global\LeqStarSectionfalse
    }
    \def\leq@section@nostar[#1]#2{%
      \global\LeqStarSectionfalse
      \gdef\leq@current@section{#2}%
      \gdef\leq@current@subsection{}%
      \leq@original@section[#1]{#2}%
    }
    % Redefinir \subsection
    \renewcommand{\subsection}{%
      \@ifstar{\leq@subsection@star}{\@dblarg\leq@subsection@nostar}%
    }
    \def\leq@subsection@star#1{%
      \global\LeqStarSectiontrue
      \gdef\leq@current@subsection{#1}%
      \leq@original@subsection{#1}%
      \global\LeqStarSectionfalse
    }
    \def\leq@subsection@nostar[#1]#2{%
      \global\LeqStarSectionfalse
      \gdef\leq@current@subsection{#2}%
      \leq@original@subsection[#1]{#2}%
    }

    % ----- Restauración de valores Globales de estilo ----------%
    % Flags
    \newif\ifInFootnote
    \InFootnotefalse
    \pretocmd{\@footnotetext}{\InFootnotetrue}{}{}
    \apptocmd{\@footnotetext}{\InFootnotefalse}{}{}
    % Profundidad de listas (soporta anidadas)
    \newcount\MyListDepth \MyListDepth=0
    \AtBeginEnvironment{la}{\advance\MyListDepth by 1}
    \AfterEndEnvironment{la}{\advance\MyListDepth by -1}
    \AtBeginEnvironment{lnum}{\advance\MyListDepth by 1}
    \AfterEndEnvironment{lnum}{\advance\MyListDepth by -1}
    \AtBeginEnvironment{itemize}{\advance\MyListDepth by 1}
    \AfterEndEnvironment{itemize}{\advance\MyListDepth by -1}
    \AtBeginEnvironment{enumerate}{\advance\MyListDepth by 1}
    \AfterEndEnvironment{enumerate}{\advance\MyListDepth by -1}
    % Restauración diferida: hazlo al INICIO del próximo párrafo real
    \newif\if@pendingRestore
    \@pendingRestorefalse
    \newcommand{\RequestRestore}{\global\@pendingRestoretrue}
    \everypar{%
      \if@pendingRestore
        \global\@pendingRestorefalse
        \ifInFootnote\else
          \ifnum\MyListDepth=0
            \setlength{\parskip}{\GlobalParskip}%
            \setstretch{\GlobalBaseLineStretch}%
          \fi
        \fi
      \fi
    }
    % Al final de bloques:
    \AfterEndEnvironment{equation}{\RequestRestore}
    \AfterEndEnvironment{equation}{\RequestRestore}
    \AfterEndEnvironment{la}{\RequestRestore}
    \AfterEndEnvironment{lnum}{\RequestRestore}
    % Al final de macros:
    \apptocmd{\eqblock}{\RequestRestore}{}{}
    \apptocmd{\imagenfit}{\RequestRestore}{}{}
    
\makeatother

% ===================== ToC propio con 4 niveles (único bloque) =====================
\makeatletter

% --- Escritores: añadimos una línea al .stc en cada cabecera numerada ---
% Guardamos las versiones actuales para llamar al comportamiento normal
\let\MyTOC@orig@section\section
\let\MyTOC@orig@subsection\subsection
\let\MyTOC@orig@subsubsection\subsubsection
\let\MyTOC@orig@paragraph\paragraph

% SECTION
\renewcommand{\section}{%
  \@ifstar{\MyTOC@section@star}{\@dblarg\MyTOC@section@nostar}%
}
\def\MyTOC@section@star#1{%
  \MyTOC@orig@section{#1}% sin entrada al .stc
}
\def\MyTOC@section@nostar[#1]#2{%
  \MyTOC@orig@section[{#1}]{#2}% comportamiento normal (escribe al .toc)
  % Escribimos también nuestra línea al .stc (texto con \numberline)
  \addcontentsline{stc}{section}{\protect\numberline{\thesection}#2}%
}

% SUBSECTION
\renewcommand{\subsection}{%
  \@ifstar{\MyTOC@subsection@star}{\@dblarg\MyTOC@subsection@nostar}%
}
\def\MyTOC@subsection@star#1{%
  \MyTOC@orig@subsection{#1}%
}
\def\MyTOC@subsection@nostar[#1]#2{%
  \MyTOC@orig@subsection[{#1}]{#2}%
  \addcontentsline{stc}{subsection}{\protect\numberline{\thesubsection}#2}%
}

% SUBSUBSECTION
\renewcommand{\subsubsection}{%
  \@ifstar{\MyTOC@subsubsection@star}{\@dblarg\MyTOC@subsubsection@nostar}%
}
\def\MyTOC@subsubsection@star#1{%
  \MyTOC@orig@subsubsection{#1}%
}
\def\MyTOC@subsubsection@nostar[#1]#2{%
  \MyTOC@orig@subsubsection[{#1}]{#2}%
  \addcontentsline{stc}{subsubsection}{\protect\numberline{\thesubsubsection}#2}%
}

% PARAGRAPH
\renewcommand{\paragraph}{%
  \@ifstar{\MyTOC@paragraph@star}{\@dblarg\MyTOC@paragraph@nostar}%
}
\def\MyTOC@paragraph@star#1{%
  \MyTOC@orig@paragraph{#1}%
}
\def\MyTOC@paragraph@nostar[#1]#2{%
  \MyTOC@orig@paragraph[{#1}]{#2}%
  % Si no numeras los \paragraph, cambia \theparagraph por texto plano sin \numberline
  \addcontentsline{stc}{paragraph}{\protect\numberline{\theparagraph}#2}%
}

% --- Impresor del índice propio (lee el .stc) ---
\newcommand{\MyTOC}{%
  \par\bigskip
  {\centering\bfseries\Large Índice de contenido\par}%
  \medskip
  \begingroup
    \parindent=0pt\parskip=2pt
    \@starttoc{stc}%
  \endgroup
}

\makeatother
% ===================== fin ToC propio =====================


%%%%%%%%%%%%%%


% --- Datos del tema ---
\title{Política Presupuestaria (II): Grandes cifras de los Presupuestos Generales del Estado\\
\large Clasificación económica, órganica y por política de gasto. Visión consolidada de los PGE}
\author{Marcelo Ortega}
\date{Septiembre 2026}

\begin{document}
\maketitle
\newpage

% --- Página de resumen y modificaciones ---
\puntosclave{ %% Escribir puntos clave Apuntes CECO
     
    }
\vspace{1cm}

\modificaciones{
    
    }

\newpage
\MyTOC
\newpage

\mylistofequations
\clearpage

\MyListOfFigures
\clearpage


\pagenumbering{arabic}
\setcounter{page}{1}


\section{Introducción}



\subsection{Contextualización}


\subsubsection{Relevancia}




\subsubsection{Problemática}



\subsection{Estructura de la exposición}

Obliga a una cierta actualización cada cierto tiempo. 

La aplicación de las reglas fiscales ha facilitado una contención del gasto público, cierta estabilidad en su estructura y el mantenimiento de su magnitud como porcentaje del PIB.\footnote{%
    La estabilidad del gasto público como porcentje del PIB no se da en el caso de los gastos relacionados con el envejecimiento de la población (pensiones, sanidad y atención a la dependencia). En este caso, la tendencia es creciente. en ausencia de medidas que la corrijan.
} Sin embargo, la crisis de la Covid-19 y la provocada por la invasión de Ucrania han dejado en suspenso las reglas fiscales durante unos años y han desatado un importante aumento del gasto sin saber bien qué parte será permanente y cuál transitoria. Del mismo modo, por el lado de los ingresos, las autoridades han aplicado rebajas de urgencia en algunos impuestos, a la vez que han ideado recargos en otros. Todo va en contra de la deseable estabilidad de las cifras a lo largo del tiempo. Es de esperar que a partir de los PGE 2023 este problema se atenúe.

No hay que perder de vista que el objeto de este tema es saber cuánto y en qué se gasta la Administración Central los recursos que recauda, y la procedencia de éstos. También se trata de saber a través de qué subsectores lleva a cabo la Administración su actividad. Por tanto, es algo sobre lo que cualquier economista español debe tener, al menos, una idea aproximada. En esto consisten las grandes cifras de las que habla el título. Por tanto, el tema no va referido a los PGE de un determinado año, sino a los rasgos más o menos permanentes que encontramos en los PGE. 

Será preciso memorizar algunas grandes cifras, pero no es necesario hacerlo con exactitud. En bastantes ocasiones se proporcionan las magnitudes en porcentaje del PIB porque dan una buena idea de la importancia de las mismas. además de ser más fácilmente memorizables y de mostrar una cierta estabilidad a lo largo del tiempo. Así. por ejemplo, si el gasto total del Estado asciende al 25,6 por ciento del PIB. el opositor puede decir que el gasto total --se sitúa en torno al 25 por ciento del PIB ... Además. las cifras que se citan son porcentajes del PI B nominal previsto en el cuadro macroeconómico de los Presupuestos Generales del Estado para 2023. Ciertamente, las cantidades de gasto ejecutadas, las de ingresos recaudados y el PlB registrado a final de año no coincidirán con lo previsto. Razón de más para que sea válido dar cifras aproximadas. En todo caso, en el tema se proporciona un número grande de cifras. con objeto de que el opositor tenga la mayor información (básica) posible y realice la selección que le parezca conveniente. 

Las cifras en valor monetario de expresan en miles de millones de euros para facilitar su memorización (en los PGE figuran en millones de euros). Dado que la principal (y, prácticamente, única) fuente bibliográfica de este tema es el Informe Económico y Financiero de los PGE, si el opositor lo ve conveniente, puede sustituir fácilmente las cifras que desee por las del último presupuesto aprobado, consultando el último Informe en internet en la página de la Dirección General de Presupuestos. También puede consultar el último proyecto ele presupuesto presentado, porque las diferencias con las que finalmente se aprneban suelen ser pequeí'ias. Para ello, en cada apartado ele este tema se consigna la tabla y la página el Informe Económico y Financiero de los PGE 2023 de la que proceden los datos. Dado que el Infonne siempre tiene la misma estructura, será fácil encontrar los datos actualizados.

A pesar de que no se trata de un tema sobre los PGE de un año en concreto (es decir. se trata de un teme «estructural»), el opositor puede considerar incluir algún comentario valorativo del (último) presupuesto aprobado, emitido por alguna institución con acreditada autoridad técnica para hacerlo. Tal sería el caso de la A IReF y de FEDEA. Hay que advertir de que estos infom1es se refieren al proyecto de presupuestos, no al presupuesto aprobado. En el Anexo I se incluye un comentario sobre aspectos contenidos en los inforn1es de estas instituciones para los PGE 2023. No obstante, el objeto del tema sigue sie1ído comentar las grandes cifras del presupuesto, más que una valoración crítica más o menos coyuntural.

Por lo demás, el tema pretende ser omnicomprensivo, de modo que se relacionan bastantes entidades cuyos presupuestos son, por comparación, poco importantes desde el punto de vista cuantitativo. Pero es legítimo que el opositor se refiera a ellas (aun sin dar cifras) pues haciéndolo se pone de manifiesto el amplio abanico de actuaciones de la Administración Central en España, lo cual es también objeto del tema.


--------------------------------- -> Introducción

El objeto de este tema es saber cuánto y en qué gasta I¡¡ Adm inistración Central sus recursos. y de dónde proceden éstos. Los datos que se proporcionarán serán de carácter aproximativo. basados en cifras de los Presupuestos Generales del Estado de 2023. Por lo tanto. se trata de importes en términos de contabi lidad presupuestaria. no de contabilidad nacional, como es el caso, por ej emplo. de los datos del Tema 23 de la Parte B del 4° ��jcrcicio. sobre las finanzas de las AAPP. Se trata, además. de importes presupuestados, no de importes ej ecutados. Efectivamente, hay diferencias entre los ingresos y gastos presupuestados y lo que tinalmente se ��jccuta o se recauda.

Los flujos presupuestarios son muchos y cambian todos los años: de ahí que nos centremos en los grandes agregados y en su importancia relativa con respecto al total y con respecto al PJB, aspectos en los que las magnitudes se mantienen a lo largo del tiempo en niveles similares. Sin embargo, esto último deja de ser cierto en el período 2022-2026 por la inclusión en los PGE, tanto por la parte de ingresos como la de gastos, del plan Next Generation EU, que cuenta con dos instrumentos principales: el Mecanismo de Recuperación y Resiliencia (MRR). que fi nancia las actuaciones del Plan de Recuperación, Transformación y Resiliercia (PRTR). y la inic iativa REACT-EU.\footnote{%
    De hecho, los programas presupuestarios correspondientes al PRTR (y al REACT-EU) tienen un código alfanumérico que los distingue de los on-os.
} Así. en el presupuesto de gastos del Estado se consignan aquellos ligados al PRTR y a la iniciativa REACT-EU2 (aunque en el Informe Económico y Financiero ambos se engloban como PRTR; al igual que haremos en este tema en aras de agilidad expositiva). Y. en el caso de los ingresos. el presupuesto se refi ere a "Fondos Next Generation-EU'·. Dado que se trata de algo transitorio, nos referiremos preferentemente al presupuesto sin PRTR, en el caso de los gastos. y al presupuesto sin fo ndos Next Generation -EU, en el caso de los ingresos. Pero en ciertas ocasiones, daremos el dato incluyendo a ambos.\footnote{%
    Por la decisión de Eurostat de aplicar el principio de neutralidad financiera, los flujos de ingresos y gasros Next Generation-El J no ti.:ncn impacto en d.Sficit cn 1érminos de contabilidad naci onal. Es decir, todo ingreso comporta un gasto equivalente. y viceversa. aunque no se produzcan en el mismo afio. Pcro se trata d.:- una norma de comabilidad nacional. no de contabilidad presupuestaria.
}

Por último, es preciso advertir que los datos presupuestarios en porcentaje del PIB son el resul tado de dividir importes calculados con distinta metodología: metodología de contabilidad presupuestaria en el numerador (gasto o ingreso), y metodología de coniabilidad nacional en el denominador (PIB ). No obstante. es una medición convencional de las magnitudes presupuestarias. También es conveniente advertir de que la similitud de las ratios sobre PlB se quiebra en el período 2020-202 1, por la crisis de la Covid- 19 que elevó el gasto y hundió el PIB y los ingresos.

Empezaremos el tema con los presupuestos del Estado. los Organismos Autónomos, la Seguridad Social y el resto de entidades del sector público adm inistrativo con presupuesto limitativo. A continuación. seguiremos con los presupuestos con carácter estimativo del sector público empresarial, fundacional y de otro tipo de entidades. Finalmente, haremos una visión de conjunto de la acción de la Administración central con los PGE consolidados.














\clearpage
\section{Contabilidad presupuestaria: Marco jurídico-conceptual}
\puntosclave{
    
}

Antes de entrar en el detalle de cómo se aplican las tres clasificaciones a cada entidad del sector público estatal, necesitamos resolver una pregunta previa que rara vez se plantea con la profundidad que merece: ¿qué es, exactamente, clasificar un presupuesto, y por qué hemos terminado clasificándolo de estas tres formas concretas y no de otras posibles? Este bloque responde a esa pregunta en tres tiempos: primero, un recorrido histórico centrado no en la institución presupuestaria en sí —ya trabajada en el Tema 21— sino específicamente en la técnica de clasificación, es decir, en cuándo y por qué apareció cada uno de los tres criterios que hoy damos por evidentes; segundo, el fundamento teórico de cada clasificación y la justificación, desde la propia teoría y no solo desde la norma positiva, de por qué se aplican únicamente a determinadas entidades; y tercero, una precisión sobre el alcance territorial de todo lo que sigue. Empiezo por la historia.

\subsection{Desarrollo histórico de la forma de presentación de los presupuestos del sector público administrativo}
Ley de Contabilidad y Administración Pública de 1850 estableció que cada ministerio debía elaborar su propio presupuesto. Conviene ahora leer ese mismo dato desde una óptica distinta. 

Esa decisión, tomada por razones prácticas de gestión es el nacimiento de la clasificación orgánica española: la primera de las tres clasificaciones en aparecer, y la que menos elaboración teórica exigió, porque se limitaba a preguntar algo que la propia estructura administrativa del Estado ya respondía por sí sola: ¿quién gasta?

La clasificación económica: una construcción más lenta, ligada a la maduración de la propia contabilidad pública
La clasificación económica —por capítulos, según cómo se gasta o se ingresa— tiene un origen mucho menos preciso y mucho más gradual.

Nace de la progresiva sofisticación de la propia contabilidad pública española a lo largo de los siglos XIX y XX, a medida que la Hacienda necesitaba distinguir, dentro del gasto de cada ministerio, qué proporción correspondía a personal, qué proporción a material y suministros, qué proporción a inversiones. Una necesidad de gestión interna que fue cristalizando en estructuras de capítulos cada vez más formalizadas.

### La clasificación funcional: la más joven de las tres, y la única con origen internacional documentado
La clasificación funcional —en qué se gasta, por programas— tiene, en cambio, un origen mucho más reciente y, esto sí, perfectamente documentado. Procede de un esfuerzo explícito de estandarización internacional promovido por Naciones Unidas a finales de los años cincuenta. 

El Departamento de Asuntos Económicos de la ONU publicó, en 1958, A Manual for Economic and Functional Classification of Government Transactions, un manual que proponía, por primera vez de forma sistemática y con vocación de aplicación universal, dos esquemas de clasificación combinados.

Uno económico —que permite conocer de inmediato los resultados económicos de la actividad del Estado— y otro funcional —que agrupa el gasto según los fines administrativos que persigue, con independencia de qué departamento concreto lo ejecute. Este manual tuvo un impacto considerable en la presentación presupuestaria y en el desarrollo de la información sobre las transacciones del sector público en un número muy amplio de países, y su recepción no tardó en llegar a España. La revista del Instituto Nacional de Administración Pública dedicó, ya en los años inmediatamente posteriores a su publicación, una nota específica a explicar el sistema que debe seguir la clasificación de los diversos conceptos del Presupuesto de Gastos para conseguir una visión más exacta de los efectos económicos de la actuación del Estado. 

El propio manual de la ONU reconocía, además, que los detalles de los esquemas de clasificación necesitarán modificarse para reflejar la estructura institucional y el alcance de la responsabilidad gubernamental predominante en cada país. Una advertencia que explica por qué la clasificación funcional española del Estado (programas, políticas de gasto, áreas de gasto) y la de la Seguridad Social (Prestaciones Económicas, Asistencia Sanitaria, Servicios Sociales, Servicios Comunes) no comparten la misma taxonomía. Cada país, y dentro de cada país cada tipo de entidad, adaptó el esquema internacional a su propia estructura institucional, en lugar de importar una plantilla única y rígida. 

Este mismo impulso de estandarización internacional se replicó después, de forma regional. En América Latina, la CEPAL convocó en Santiago de Chile en 1962 y en San José de Costa Rica en 1963 sendos Seminarios de Clasificación y Administración Presupuestarias, de los que surgió, en 1964, un Esquema Uniforme de Clasificación de las Cuentas del Sector Público adaptado a los países de América Latina.

### La presupuestación por programas: cuando la clasificación funcional deja de ser solo una forma de presentar y se convierte en una técnica de gestión

Conviene conectarlo con el debate de la presupuestación por programas. España adoptó desde los años 80, no es una simple evolución técnica de la clasificación funcional de 1958. Es, específicamente, la llegada a España del movimiento internacional de presupuesto por programas (Planning-Programming-Budgeting System, PPBS) que ya analizamos al hablar del principio de eficiencia, con Robert McNamara, Lyndon Johnson y la crítica incrementalista de WILDAVSKY. Es decir, la clasificación funcional española de hoy hereda dos genealogías distintas que confluyen en el mismo punto. La genealogía puramente clasificatoria de 1958 (agrupar el gasto por fin, con independencia de quién lo ejecute) y la genealogía de gestión por resultados del PPBS de los años sesenta (no solo clasificar el gasto por fin, sino evaluarlo por el grado de cumplimiento de ese fin mediante indicadores).\footnote{%
    La propia AIReF ha señalado que esta técnica tiene en España "una aplicación más formal que real" no es, por tanto, una crítica aislada al caso español: es la reaparición, sesenta años después y en otro país, de la misma tensión entre ambición racionalizadora y resistencia práctica de la negociación política que Wildavsky ya documentó para el propio fracaso del PPBS estadounidense en 1971.
}


\subsection{El porqué y el quién de cada clasificación}
Este recorrido histórico deja ya planteada la siguiente pregunta que abordaremos. Si las tres clasificaciones tienen orígenes tan distintos, ¿comparten realmente una misma lógica de fondo, o son, en realidad, tres respuestas a tres preguntas genuinamente distintas que conviene tratar por separado?

La tentación más habitual es tratarlas variantes de una misma idea: formas distintas de ordenar los mismos números. Es un error: cada una responde a una pregunta analítica distinta, nacida de una tradición intelectual distinta, con defensores y críticos propios.

La orgánica pregunta quién es responsable del gasto. La económica pregunta qué naturaleza económica tiene ese gasto. La funcional pregunta para qué fin se realiza. Tres preguntas que, aplicadas al mismo euro gastado, pueden dar tres respuestas completamente independientes: el mismo euro puede estar bajo la responsabilidad del Ministerio de Sanidad (orgánica), tener naturaleza de transferencia corriente (económica), y perseguir el fin de protección social si se trata de una prestación por dependencia (funcional). 

\subsubsection{Clasificación orgánica: la pregunta de la responsabilidad política}
Su fundamento teórico no es, en realidad, presupuestario, sino constitucional.

Hunde sus raíces en la misma doctrina de responsabilidad ministerial que tocamos al hablar del origen inglés del poder de la bolsa ([\authorfont{Ver Tema 4.B.22}]). La tradición constitucional de Westminster estableció, desde el siglo XVIII, que cada ministro responde individualmente ante el Parlamento de la gestión de su propio departamento: principio de responsabilidad ministerial individualizada, distinta a la naturaleza colegiada de nuestro ordenamiento. 

En este linea, la clasificación orgánica es la traducción presupuestaria exacta de esa doctrina: si un ministro debe responder políticamente de su gestión, el presupuesto tiene que poder decirnos qué gasto concreto cae bajo su responsabilidad. Es, por esta razón, la clasificación más antigua de las tres, precisamente porque la estructura ministerial del Estado ya preexistía a cualquier reflexión sobre cómo clasificar el presupuesto: la orgánica no tuvo que inventar su objeto de análisis, se limitó a heredarlo de la propia arquitectura del poder ejecutivo. La crítica clásica, con especial fuerza desde la tradición de presupuestación por programas, es que, precisamente por organizarse en torno a la estructura administrativa y no en torno al fin perseguido, tiende a fragmentar artificialmente políticas públicas que, en la realidad, son transversales a varios departamentos. 

Pensemos la lucha contra el cambio climático. Evidentemente, la transición ecológica involucra gasto del Ministerio para la Transición Ecológica (renovables, gestión hídrica), del Ministerio de Transportes (movilidad sostenible), del Ministerio de Agricultura (adaptación del sector primario) y del Ministerio de Industria (descarbonización industrial). Una clasificación puramente orgánica jamás permitiría responder a la pregunta ¿cuánto gasta España en lucha contra el cambio climático? Cada ministerio informaría de su parte, pero nadie las sumaría de forma natural dentro de ese esquema. Exactamente esta insuficiencia motivó, como vimos, la creación de la clasificación funcional en el manual de la ONU de 1958.

Desde una óptica internacional, el peso relativo que cada tradición presupuestaria concede a la orgánica frente a la funcional varía notablemente. Los países de tradición continental europea (Francia, España) han mantenido históricamente una fuerte primacía de la clasificación orgánica como eje vertebrador del documento presupuestario, coherente con su tradición de responsabilidad ministerial estricta. Los países que más lejos han llevado las reformas de Nueva Gestión Pública (Nueva Zelanda, Australia y Reino Unido) han invertido esa primacía, organizando sus documentos presupuestarios centrales en torno a resultados y funciones (\textit{outputs}, \textit{outcomes}), relegando la estructura departamental a un segundo plano informativo.

\subsubsection{Clasificación económica: la pregunta de la macroeconomía, y su genealogía hasta un premio Nobel}
Esta es, con diferencia, la clasificación con la raíz teórica más profunda y mejor documentada de las tres. 

Su origen no está en la administración pública, sino en el nacimiento mismo de la contabilidad nacional como disciplina económica. STONE, junto con MEADE, desarrolló, durante la Segunda Guerra Mundial y por instigación directa de KEYNES, el primer sistema de cuentas nacionales del Reino Unido. Publicado en 1941 dentro de un informe gubernamental que se titulaba \textit{An Analysis of the Sources of War Finance}: la necesidad bélica de saber con precisión de dónde salía cada libra que financiaba el esfuerzo de guerra británico es el origen práctico e inmediato de toda la contabilidad nacional moderna. 

STONE fundaría después, en 1945, el \textbf{Departamento de Economía Aplicada de Cambridge}, y publicaría en 1953 y 1968 sucesivas versiones de \textit{A System of National Accounts}, que se convertirían en la base directa del Sistema de Cuentas Nacionales de Naciones Unidas (SEC-2010), mereciendo por ello el Premio Nobel de Economía de 1984.\footnote{%
    Una de las tradiciones heterodoxas más influyentes en el ámbito financiero surge, precisamente, del Departamento de Economía Aplicada de Cambridge: el enfoque stock-flujo, de GODLEY. La tradición metodológica, más que teórica, de este Departamento, ha sido muy influyente en la literatura económica reciente y las nuevas corrientes económicas.
}

Pero, ¿por qué este linaje académico en el fundamento de la clasificación económica? Por una sencilla razón, para que la contabilidad nacional de STONE pudiera integrar de forma coherente al sector público dentro del conjunto de la economía, necesitaba disponer del gasto público desglosado en categorías económicamente significativas y no meramente administrativas. Desde la óptica de STONE, no importa si un gasto lo ejecuta el Ministerio de Educación o el de Sanidad. Lo que importa, y mucho, es si ese gasto es consumo (con un efecto multiplicador distinto) o inversión (con un efecto sobre el stock de capital productivo), si es una transferencia (que redistribuye renta sin generar directamente producción) o una compra de bienes y servicios (que sí computa directamente en el PIB por el lado del gasto). 

La clasificación económica presupuestaria es la puerta de entrada que permite conectar el presupuesto del Estado con el aparato analítico macroeconómico completo que STONE diseñó. Y es, por ello, la única de las tres clasificaciones cuya lógica de capítulos (personal, bienes corrientes, transferencias, inversiones) se corresponde con las categorías económicas que trabajamos al estudiar las finanzas de las AAPP ([\authorfont{Ver Tema 4.B.23}]).

No obstante, precisamente por estar diseñada para el análisis macroeconómico agregado, la clasificación económica resulta poco informativa para la gestión cotidiana de un programa concreto. Saber que un ministerio ha gastado \textit{tanto en transferencias} no dice nada, por sí solo, sobre si ese gasto ha sido eficaz o ineficaz para el fin que perseguía.

\subsubsection{Clasificación funcional: la pregunta de las prioridades, entre la transparencia democrática y el problema de las fronteras arbitrarias}
Su origen, como ya dijimos, es el manual de la ONU de 1958 y su conexión con el PPBS.

No obstante, su fundamento teórico de fondo es de naturaleza democrática. ¿Por qué? Porque la clasificación funcional permite al ciudadano y al Parlamento conocer las prioridades reales del Estado: cuánto se dedica a sanidad, a educación, a defensa; con independencia de la arquitectura administrativa concreta que las ejecute, que puede cambiar de nombre y de estructura en cada reorganización ministerial sin que las prioridades de fondo varíen. Es, en este sentido, la clasificación que mejor conecta con la idea de que el presupuesto es un \textit{contrato entre la sociedad y su Gobierno}: \textbf{el ciudadano no vota por cuánto gasta el Ministerio de Hacienda en consumo corriente; vota por cuánto se dedica a sanidad o a pensiones, con independencia de qué departamento concreto lo gestione.}

Como las demás, no está exenta de críticas. De hecho, su crítica interna se dirige al hecho de que por agrupar el gasto según fines y no según estructura administrativa objetivamente verificable, se enfrenta a un problema de fronteras arbitrarias. Es decir, un mismo gasto puede, con la misma legitimidad, adscribirse a más de una función.\footnote{%
    Una crítica más general, presentada en [\authorfont{Ver Tema 4.B.21}], sería la de WILDAVSKY.
}

Por ejemplo, una inversión en infraestructura ferroviaria de alta velocidad, ¿es transporte (función económica), cohesión territorial (protección social en sentido amplio), o actuación de carácter económico (fomento de la actividad productiva)? La clasificación por políticas de gasto exige tomar decisiones de encuadre que, inevitablemente, contienen un margen de discrecionalidad mayor.

### La justificación teórica de la línea divisoria: por qué solo el presupuesto limitativo necesita clasificación

Llegamos ahora a la pregunta que da sentido a toda la restricción de alcance que fijamos para el tema: ¿por qué solo las entidades con presupuesto limitativo usan estas tres clasificaciones, y no todas las entidades del sector público? 

La razón es una consecuencia lógica de la naturaleza del presupuesto limitativo. Recordemos que el principio de especificación es el límite, cuantitativo y cualitativo, que el Parlamento impone al Ejecutivo sobre el presupuesto. Pero para que este límite pueda enunciarse siquiera, necesita un lenguaje capaz de decir \textit{hasta aquí, y no más, en esta categoría concreta}. Ese lenguaje es, exactamente, la clasificación presupuestaria. Sin clasificación orgánica, no podríamos decir el Ministerio de Sanidad no puede gastar más de X. De la misma forma, sin clasificación económica, no podríamos decir el Capítulo I de personal no puede superar Y. Y, por último, sin clasificación funcional, no podríamos vincular un crédito a un programa concreto y verificar si se ha desviado de su fin. 

La clasificación es la infraestructura técnica que hace posible la vinculación de créditos, y la vinculación de créditos es, a su vez, la que hace posible el principio de especificación.

De aquí se sigue directamente la razón por la que las entidades con presupuesto estimativo no necesitan (y, de hecho, no podrían usar de forma coherente) este mismo aparato clasificatorio. Su presupuesto no es un límite legal que el Parlamento impone y cuyo respeto haya que verificar, sino una previsión de ingresos y gastos, más próxima a la cuenta de resultados de cualquier empresa, que puede superarse sin que ello constituya, por sí solo, ninguna irregularidad jurídica. Si no hay límite legal que vigilar, no hace falta el lenguaje técnico diseñado específicamente para vigilarlo.\footnote{%
    Por ejemplo, ENAIRE o el ICO presentan un presupuesto de explotación y de capital, el lenguaje contable de cualquier empresa, y no un presupuesto por capítulos I a IX con programas y memorias de objetivos. Sencillamente, no tienen nada equivalente a la especificación parlamentaria que ese lenguaje fue diseñado para servir.
}

Esta misma lógica explica, por extensión, por qué las Comunidades Autónomas y las Entidades Locales no consolidan, sin embargo, dentro de los PGE en sentido estricto. Como sabemos, estas Instituciones sí tienen presupuesto limitativo y, por tanto, necesitan y usan un aparato clasificatorio análogo al que hemos descrito. Pero, ¿por qué no se ve consolidado? Como vemos, no es que su presupuesto no esté clasificado, sino que cada una aprueba y ejecuta su propia especificación parlamentaria, ante su propio Parlamento autonómico o su propio Pleno municipal, de modo que no tiene sentido agregar sus créditos dentro de la especificación que las Cortes Generales imponen, específicamente, al Estado y a sus entidades dependientes. 

La frontera entre \textit{lo que entra} en los PGE y \textit{lo que no} no es, por tanto, una cuestión de tamaño ni de importancia relativa. Es una cuestión de ante qué Parlamento concreto responde cada especificación presupuestaria.\footnote{%
    De hehco, es una pregunta análoga a la de responsabilidad política que, no por casualidad, sostenía el fundamento teórico de la propia clasificación orgánica.
} Las tres clasificaciones, y la propia frontera de qué entidades las usan, terminan remitiendo a la misma pregunta constitucional de fondo: quién debe rendir cuentas, ante quién, y con qué grado de detalle.

\subsection{La extensión territorial: los PGE y los niveles inferiores de Gobierno}
Todo el desarrollo de este tema se centra, como ya anunciamos, en los Presupuestos Generales del Estado en sentido estricto: el Estado y sus entidades dependientes con presupuesto limitativo. 

Pero el propio epígrafe anterior ya nos ha dado la clave para entender qué ocurre en los niveles inferiores de Gobierno. Si la clasificación presupuestaria es la infraestructura técnica necesaria para que el principio de especificación tenga sentido, y ese principio se predica de cualquier ente con presupuesto limitativo y no solo del Estado, entonces las Comunidades Autónomas y las Entidades Locales necesitan un aparato clasificatorio análogo. 

La pregunta que queda por resolver no es si existe esa réplica, sino cuánto se parece a la del Estado, y por qué.

\subsubsection{Nivel local: una armonización mandatada por ley, no una simple coincidencia de estilo}
Para las Entidades Locales, la respuesta está legalmente mandatada. El art. $167.1$ del Texto Refundido de la Ley Reguladora de las Haciendas Locales (RD Legislativo 2/2004) encomienda expresamente al Ministerio de Hacienda el establecimiento, con carácter general, de la estructura de los presupuestos de las entidades locales, teniendo en cuenta la naturaleza económica de los ingresos y de los gastos, así como las finalidades u objetivos que con estos últimos se propongan conseguir.

La misma tríada orgánica/económica/funcional que hemos visto, impuesta por ley a los más de ocho mil municipios y provincias españolas. Y, de forma todavía más reveladora, la disposición adicional sexta de esa misma ley obliga al Ministerio a modificar esa estructura con objeto de adaptarlos a los establecidos para el sector público estatal en cada momento, una cláusula de convergencia automática y permanente que garantiza que la clasificación local nunca pueda desviarse de forma duradera de la que rige para el Estado.

La norma actualmente vigente que desarrolla este mandato es la Orden EHA/3565/2008, que sustituyó a su vez a: la Orden de 20 de septiembre de 1989, la primera estructura presupuestaria local obligatoria, aplicable desde 1992; y a la Orden EHA/1645/2004, que había adaptado la estructura al cambio de nomenclatura que el Estado había introducido, la sustitución de la vieja clasificación funcional por la actual clasificación por programas, vinculada a áreas y políticas de gasto. 

Conviene señalar que la Orden EHA/3565/2008 hace obligatorias, como mínimo, las clasificaciones por programas y económica, pero deja la clasificación orgánica como opcional para cada Entidad Local, en función de su propia complejidad organizativa. Un ayuntamiento pequeño, con una estructura administrativa mínima, puede prescindir razonablemente de ella sin perder información, mientras que una gran Diputación Provincial, con múltiples Organismos Autónomos dependientes, debería aplicarla para ganar detalle. Además, la propia norma reconoce que esta convergencia no es total, con el fin de adaptarse a las diferencias estructurales que existen entre Estado y Entidades Locales.\footnote{%
    Por ejemplo, dentro de la clasificación por programas, seria estúpido asimilar todos los programas del área de gasto en \textit{Servicios Públicos Básicos}. Como es normal, tiene un contenido sustancialmente distinto (abastecimiento de agua, recogida de residuos, alumbrado público, policía local) precisamente porque las funciones que un municipio presta a sus vecinos son distintas de las que presta el Estado a la nación en su conjunto.
}

\subsubsection{Nivel autonómico: mayor autonomía normativa, la misma lógica de fondo}
En contraste con lo anterior, las Comunidades Autónomas no están sujetas a armonización ministerial estatal. Cada una regula la suya mediante su propia Ley de Hacienda Pública autonómica, en ejercicio de la autonomía financiera. 

Sin embargo, esta mayor autonomía normativa no ha generado una fragmentación real de criterios. Todas las Comunidades Autónomas de régimen común han adoptado, de forma voluntaria y convergente, un esquema de clasificación orgánica/económica/por programas equivalente al estatal, precisamente porque la necesidad de que sus cuentas puedan consolidarse dentro del agregado de contabilidad nacional de las AAPP exige, como condición técnica previa, un mínimo de comparabilidad clasificatoria entre todos los subsectores. Sin esa comparabilidad, ni la IGAE podría construir la Central de Información Económico-Financiera de las Administraciones Públicas que agrega los datos de todos los niveles de gobierno, ni sería posible verificar, subsector por subsector, el cumplimiento de las reglas de déficit, deuda y gasto de la LOEP. La armonización clasificatoria territorial no es, por tanto, un capricho de uniformidad burocrática, sino la precondición técnica de todo el edificio de gobernanza fiscal multinivel.

\begin{specialblock}{\textbf{Territorios Forales:} Excepción}
    Esta convergencia general tiene su excepción natural en los territorios forales. El País Vasco y Navarra, al gestionar su propia Hacienda con un grado de autonomía institucional muy superior al del resto de Comunidades Autónomas, mantienen también sus propios esquemas de clasificación presupuestaria, sin la misma sujeción a la convergencia automática que la disposición adicional sexta de la LRHL impone al nivel local. 
    
    Es, una vez más, la misma asimetría de fondo reapareciendo aquí, en el terreno aparentemente más técnico y menos político de todos los que hemos recorrido en este tema. Ni siquiera la clasificación presupuestaria, con su vocación de tecnicismo neutral, escapa por completo a la huella de la historia constitucional española que atraviesa, como hemos comprobado una y otra vez a lo largo de esta serie, prácticamente cualquier institución de la Hacienda Pública española.
\end{specialblock}



\clearpage
\section{Contabilidad presupuestaria: Impacto material}
\puntosclave{
    La clasificación orgánica divide el Presupuesto del Estado en 38 secciones. 
    
    Los principales agregados del gasto son: gasto total. gasto de operaciones no financieras (también llamado «gasto no financiero») . y lim ite de gasto no financiero (conocido como . . techo de gasto')
    
    El gasto total del Estado ronda los 356mm de euros (un 25% del PIB ). El PRTR le ariade un 8% más. El gasto de operac iones no financieras es algo más del 16%, del PfB; y el techo de gasto es el 12,5% del PIB. 
    
    Si a las cifras anteriores les restamos los gastos que hay que hacer obligatoriamente (intereses de la deuda, transferencias a la Segu1idad Social y al SEPE para pensiones y desempleo. apo1tación al presupuesto de la UE etc.), la cantidad disponible para los ministerios es sólo de 54mm. Si restamos también los gastos de personal y de compras de bienes y servicios, la suma se queda en 28.9mm (un 2. 1 % del PIB y un 16,7% del techo de gasto). 
    
    Las transferencias corrientes son el capítulo de gasto más importante. con el 63% del gasto de operaciones no financieras. Revela el papel central del Estado en los PGE y en el marco presupuestario nacional. pues este capítulo recoge flujos de financiación a la Seguridad Social. al SEPE a las CCAA y a la CCLL además de la aportación al presupuesto de la UE 
    
    El capítulo VIII. Activos financieros. tiene importancia porque, mediante préstamos a empresas (públicas y privadas), el Estado fo menta la inversión productiva. 
    
    En el caso de los ingresos. destacan los ingresos impositivos. siendo los más importantes: IRPF, IV A, Impuesto sobre Sociedades y los Impuestos Especiales.
}


Cerrado el aparato teórico del Bloque 1, entramos en su aplicación. Recorremos, entidad por entidad, las tres clasificaciones que acabamos de fundamentar —orgánica, económica y funcional—, empezando por el Estado, la pieza central y de mayor peso del presupuesto consolidado, y continuando por los Organismos Autónomos, la Seguridad Social y el resto de entidades con presupuesto limitativo. En cada una, las "grandes cifras" surgirán de forma natural al hilo de cada clasificación, tal y como acordamos, y no como un bloque numérico separado. No me limito a reproducir las cifras del documento de referencia: cada una se acompaña de su porqué, de su genealogía cuando la tiene, y de los debates de fondo que la cifra desnuda no deja ver.

\subsection{Estado}

Clasificación orgánica

El Estado se estructura, orgánicamente, en 38 secciones.

Los órganos constitucionales constituyen seis secciones: Casa de S.M. el Rey, Cortes Generales, Tribunal de Cuentas, Tribunal Constitucional, Consejo de Estado y Consejo General del Poder Judicial. Una sección será ocupada por cada uno de los Ministerios de la legislatura (o periodo presupuestario, idealmente), en la actualidad, son veintidós departamentos ministeriales. Por último, hay una serie de secciones que no se encuadran en ninguna de las anteriores, formando un todo heterogéneo. En la actualidad, son nueve secciones.\footnote{%
    De diversa índole. Las actuales son: Aportaciones al Mutualismo Administrativo, Deuda Pública, Clases Pasivas, Contratación Centralizada, Otras Relaciones Financieras con Entes Territoriales, Fondos de Compensación Interterritorial, Relaciones Financieras con la Unión Europea, Fondo de Contingencia y Sistemas de Financiación de Entes Territoriales.
}

Esta última categoría revela algo importante sobre la propia arquitectura. Por ejemplo, una de éstas es la sección de Deuda Pública, que el hecho de que no dependa de ningún Ministerio concreto responde, en cierto modo, a la exigencia del art. 135.3 CE, que otorga a los créditos para satisfacer los intereses y el capital de la deuda pública una prioridad absoluta de pago. En este caso, dotarlos de una sección orgánica propia, ajena a la jerarquía ministerial es la traducción exacta de ese estatus privilegiado: una sección que no puede quedar sujeta a la misma lógica de responsabilidad política departamental que sí rige para el resto de secciones.

El reparto en términos de cuantía entre las $38$ secciones. Es muy desigual. Los Órganos Constitucionales apenas vinculan el $0,2\%$ del total. El conjunto de secciones ministeriales compromete una cuantía del $44\%$. Y, por último, el $55,8\%$ se distribuye entre secciones como Deuda Pública, Clases Pasivas, o los fondos de financiación a otras Administraciones.

Por tanto, como podemos ver una parte sustancial del Presupuesto del Estado no es gestión ministerial discrecional, sino compromisos estructurales que trascienden la propia lógica de responsabilidad ministerial individual que fundamenta, en teoría, la clasificación orgánica. Dentro del $44\%$ que sí corresponde a Ministerios, destaca de forma muy señalada el Ministerio de Inclusión, Seguridad Social y Migraciones, con el $36,29\%$ del total ministerial, que incluye las transferencias del Estado a la Seguridad Social para financiar los gastos impropios ([\authorfont{Ver Tema 4.B.23}]) y prestaciones no contributivas de gestión directa como el Ingreso Mínimo Vital.\footnote{%
    Este porcentaje no incluye, por tanto, las pensiones contributivas de jubilación, viudedad o incapacidad, que pertenecen al presupuesto de la entidad Seguridad Social, un subsector distinto.
}

Clasificación económica del gasto: la cadena de restas sucesivas, y por qué el margen de maniobra se ha ido estrechando

El gasto total del Estado en PGE23 asciende a \textcolor{red}{356,1 mil millones de euros}, el $25,6\%$ del PIB previsto. 

De ellos, el gasto de operaciones no financieras (Capítulos I a VII) suma $223,1$mm, el $16,1\%$ del PIB. Restando los créditos de programas especiales de modernización de Defensa y los gastos del sistema de financiación de CCAA y EELL ($45,1$mm), se obtiene el límite de gasto no financiero o techo de gasto: $173,1$mm, el $12,5\%$ del PIB.

Este techo no es una cifra que el Gobierno calcule libremente, sino que responde a la fórmula fijada en el art. $30.1$ LOEP: previsión de ingresos, más el objetivo de déficit aprobado, más los ajustes de contabilidad nacional, menos la financiación destinada a CCAA y EELL. Es, en cierto sentido, el punto donde toda las implicaciones de estabilidad presupuestaria se convierten en cifra concreta y operativa ([\authorfont{Ver Tema 4.B.21}]).

Del techo de gasto ($173,1$mm) hay que descontar todavía los gastos obligatorios no ministeriales. Es decir, transferencias a la Seguridad Social, al SEPE, intereses de deuda, Clases Pasivas, aportación al presupuesto de la UE, Fondo de Contingencia y mutualismo administrativo. Tras este descuento, quedan $54,4$mm ($3,9\%$ PIB) realmente disponibles para el conjunto de los Ministerios. 

Descontando también personal y bienes/servicios corrientes, gastos ineludibles a corto plazo aunque de gestión ministerial, el margen se reduce a $28,9$mm, apenas el $2,1\%$ del PIB y el $16,7\%$ del propio techo de gasto. Conviene detenerse en esta cifra final porque plantea una pregunta relevante para entender la política presupuestaria: ¿ha sido siempre tan reducido este margen, o es un fenómeno que se ha ido acentuando con el tiempo? 

La literatura comparada de finanzas públicas describe, para prácticamente todas las democracias desarrolladas, un fenómeno análogo bajo el nombre de crecimiento del gasto obligatorio o de derecho frente al gasto discrecional.\footnote{%
    En el debate presupuestario estadounidense, la distinción entre \textit{mandatory spending} (Seguridad Social, Medicare, intereses de deuda, que se pagan automáticamente por mandato legal sin necesidad de aprobación anual específica) y \textit{discretionary spending} (el resto, sujeto a aprobación anual del Congreso) es, de hecho, explícita
} El envejecimiento demográfico, que incrementa año tras año el peso estructural de las pensiones y del gasto sanitario, y el propio servicio de una deuda pública que ronda hoy el $100\%$ del PIB, son las dos fuerzas de fondo que explican por qué este margen de maniobra tiende, estructuralmente y en prácticamente todos los países comparables, a estrecharse con el paso del tiempo. Cada nuevo compromiso legal de gasto que se aprueba hoy se convierte, de facto, en gasto obligatorio del mañana, reduciendo el margen disponible para las prioridades políticas de cualquier Gobierno futuro.


Sobre la composición del gasto debemos destacar también algunas cosas. 

Las transferencias corrientes (Cap. IV) son, con diferencia, el capítulo más importante, ocupando el $39,4\%$ del gasto total, $63\%$ del gasto no financiero, reflejando el papel del Estado como redistribuidor central entre administraciones y hacia familias e instituciones. Por último. hay que decir que hasta 2026, habrá importantes transferencias corrientes a las Administraciones Te1Titoriales debidas al PRTR. Así. en 2023, se consignaron transferencias corrientes directas a Administraciones Territoriales con cargo al MRR por importe de 7.748 millones de euros.

En resumen, en 2023 los principales destinatarios de las transferencias co1Tientes del Estado, incluyendo el PRTR, son: la Seguridad Social (27,1 por ciento), las CCAA ( 1 9,6 por ciento). las familias ( 17,5 por ciento). los EELL ( 16 por ciento) y Exterior ( 13,4 por ciento). La importancia relativa de la empresa privada es, en cambio, muy reducida: un 0.5 por ciento.

El servicio de la Deuda (Cap. III) es el segundo capítulo en importancia, con un comportamiento anticíclico bien documentado. Sube en las crisis, cuando el propio endeudamiento se dispara para financiar el déficit adicional, y baja en las expansiones. Conviene advertir, que el cómputo presupuestario del gasto en intereses no coincide con el de contabilidad nacional.

Personal (Cap. I) ocupa el tercer lugar, concentrado en los ministerios de mayor plantilla (Justicia, Defensa, Interior, Educación). Junto con las compras de bienes y servicios (Capítulo II), de menor importancia cuantitativa, constituyen los gastos operativos del Estado. Los gastos del Capitulo II experimentan incrementos los años en que hay que financiar procesos electorales (Ministerio del Interior), o en los que hay que incrementar la atención humanitaria (Ministerio de Inclusión, Migraciones y Seguridad Social), o en los que el ejército debe realizar operaciones en el exterior (Defensa).

El Capítulo VII, aunque de menor cuantía que el Capítulo IV, también participa en el carácter central de aquél, pues contiene el Fondo de Compensación Interterritorial (FCI). Además, es una palanca de fomento de la inversión de empresas públicas y privadas. En este capítulo destacan las transferencias a la CNMC para financiar costes del sistema eléctrico y gasista, las transferencias para programas de investigación, las transferencias del Ministerio de Sanidad a las CCAA para mejorar las infraestructuras sanitarias y las transferencias a RENFE para la financiación de diversas actuaciones. Junto con el Capítulo VI (inversiones reales), engloba las operaciones de capital del presupuesto. Las inversiones reales tienen una importancia cualititativa menor. Sin embargo, la verdadera importancia del Estado en la Formación Bruta de Capital Fijo será puesta de manifiesto más adelante. En este capítulo destacan los gastos de los programas especiales de armamento del Ministerio de Defensa (que pueden superar el $50\%$ del total), los del Ministerio de Transpones, Movilidad y Agenda Urbana (fundamentalmente a la conservación de carreteras y la creación de nuevas infraestructuras) y los del Ministerio para la Transición Ecológica y el Reto Demográfico. espec ialmente los destinados a infraestructuras hidráulicas y a actuaciones en el medio natural).

La suma de los capítulos VI y VII asciende al $10\%$ por ciento del gasto no financiero. 

El Fondo de Contingencia de ejecución presupuestaria (Capítulo V), como dispone el art. $50$ de la LGP, es el $2\%$ del gasto no financiero, excluidos los destinados a financiar a las comunidades autónomas y entes locales en aplicación de sus respectivos sistemas de financiación. Su función es atender las necesidades de gasto (corriente y de capital) no discrecional y no previstas que puedan presentarse durante el ejericio. Conviene añadir un matiz. Su dotación legal, $2\%$ del gasto no financiero excluida la financiación de CCAA/EELL, no es una cifra elegida al azar, sino un porcentaje que pretende equilibrar dos riesgos opuestos. Por un lado, una dotación demasiado pequeña dejaría al Estado sin margen real para responder a imprevistos sin recurrir sistemáticamente a créditos extraordinarios con rango de ley. Por otro lado, una dotación demasiado grande equivaldría a introducir una bolsa de gasto discrecional adicional fuera del escrutinio detallado que la propia Ley de Presupuestos exige para el resto de partidas.

Los capítulos financieros (Capítulos VI 11 y IX) suelen pasarse por alto al analizar el presupuesto porque no computan a la hora de calcular el déficit público. Se trata de \textit{operaciones por debajo de la línea}. Sin embargo, sí computan para determinar la necesidad de endeudamiento.

La importancia económica de los préstamos del Estado (Capítulo VIII) es indudable. Aparte de la cuantía misma del capítulo, porque el Estado influye influye en el sector productivo de la economía mediante la concesión de crédito. Destacan las aportaciones al Fondo de Financiación a las CCAA, los préstamos a la Seguridad Social para su equilibrio financiero y los proyectos técnico-industriales de carácter militar. En este último caso, el Estado concede préstamos a las empresas fabricantes para dotarlas de medios económicos, con independencia de los pagos intermedios y finales a medida que avance el proyecto. Por último, cabe destacar que si el prestatario no devolviese el préstamo, esa incobrabilidad se reclasificaría como transferencia de capital, que sí computaría en déficit. Por lo tanto, es también un mecanismo que permite al Estado conceder ayuda financiera con una contabilización inicial más benigna que la de una subvención directa, siempre que exista una expectativa razonable de recuperación del principal. 

Por su parte, el importe del Capítulo IX recoge la amortización de la Deuda y los préstamos del Estado en el período, por una cuantía que alcanza casi el $30\%$ del gasto total.

Clasificación económica del ingreso: la jerarquía tributaria, y el peso simbólico del IRPF frente al IVA

Los ingresos no financieros del Estado ascienden a $171$mm ($12,3\%$ del PIB), de los cuales $144,3$mm ($84,4\%$) son ingresos impositivos por IRPF ($51,6$mm), IVA ($42,5$mm), Impuesto sobre Sociedades ($28,5$mm), Impuestos Especiales ($8,9$mm, con Hidrocarburos como partida mayor).

Conviene detenerse en algo que muestran las cifras. El IRPF supera al IVA como primera figura recaudatoria del Estado, pese a que el IVA es la segunda figura tributaria por recaudación en España tras el propio IRPF a nivel de conjunto de AAPP. La aparente contradicción se resuelve fácilmente recordando el sistema de cesión que ya trabajamos: buena parte del IVA recaudado en territorio español no llega nunca a computar como ingreso del Estado en esta clasificación, porque se transfiere a las Comunidades Autónomas dentro de su sistema de financiación, mientras que el IRPF, pese a estar igualmente cedido en un porcentaje sustancial, conserva aquí una fracción mayor bajo gestión y contabilización estatal directa. 

Al igual que en el presupuesto de gastos, los ingresos por transferen cias co1Tientes reflejan la centralidad del Estado en el funcionamiento de las AAPP. Así. la mayor parte proviene de otras Adm i nistraciones Públicas y de Organismos Autónomos. Las CCAA transfieren al Estado por el Fondo de Suficiencia (casi un 47 por ciento del capítulo IV). Por su parte. los ingresos de los EETT suponen más del 13 por ciento de este capítulo. Así mismo. un 32 por ciento de este capítulo está engrosado por fondos Mecanismo de Recuperación y Resiliencia (MRR) y del REACT EU (es decir, de los fondos Next Generation EU).

Transferencias de capital. Las transferencias de capital recogen, principalmente, los ingresos provenientes del FEDER, del MRR y del REACT EU. Suele ser un capítulo que fluctúa de año en año, puesto que los ingresos provenientes del FEDER no muestran unifonnidad a lo largo del tiempo, pues dependen del ritmo de ej ecución y de certificación de los proyectos que se financian. En el caso de los fondos Next Generation EU (MRR y REACT EU), en 2023 engrosaron el 98 por ciento del capítulo.

Ingre sos patrimoniales. Los ingresos patrimoniales son cuantitativamente poco importantes (6,2mm). Sólo cabe mencionar que entre los mismos se incluyen los dividendos y participaciones en beneficios (por ej emplo, los del Banco de España) e intereses de préstamos concedidos a Administraciones Territoriales, los ingresos de las subastas de derechos de gases de efecto invernadero y los ingresos procedentes por las subastas del espectro radioeléctrico de la telefonía 5G.

Ingresos Financieros. En 2023 se presupuestaron ingresos de 3,5 mm en el Capítulo Vlll (Activos Financieros). Son la devolución de préstamos que en su día hizo el Estado. No computan en el cálculo del saldo público Aunque en el presupuesto de ingresos no hay un Capítulo IX para registrar los ingresos por emisión de deuda del Estado, ésta varía durante el afio. La variación de la deuda del Estado está determinada por la emisión de valores y contratación de préstamos que realiza el Tesoro Público para hacer frente a las necesidades de financiación y a los vencimientos que se producen en el ejercício.

Las necesidades de financiación del Estado se derivan del déficit de caja no financiero del Estado y de la variación neta de activos financieros (la diferencia entre el capítulo V i l I de gastos y el capítulo V I I I de i ngresos) 26, incluyendo los gastos del PRTR y los ingresos de los fondos Next Generation-EU. Estas necesidades se cubren con emisiones de deuda del Estado y, en su caso, con financiación neta captada y no aplicada en ejercicios anteriores. Adicionalmente, el Tesoro debe refinanciar deuda emitida en ejercicios anteriores. La suma de las necesidades de financiación y las refinanciaciones del ejercicio determinará las emisiones brutas del Tesoro.

La emisión neta del Tesoro Público prevista para 2023 se situó en 70.000 mil lones. Añadiéndole la refinanciación de l a deuda con vencimiento en el año, la emisión bruta prevista para 2023 se situó en 256.930 millones.

\paragraph{Saldo del Estado: dos conceptos, dos propósitos distintos}

El saldo presupuestario, diferencia entre ingresos y gastos de operaciones no financieras, en términos de derechos y obligaciones reconocidas, no coincide con la capacidad o necesidad de financiación en contabilidad nacional ([\authorfont{Ver Tema 4.B.23}]).

Esto se debe esencialmente a divergencias en el principio contable utilizado. La contabilidad nacional usa el criterio de derechos y obligaciones reconocidos, siguiendo precisamente la lógica de devengo puro de STONE. Por otro lado, el saldo de caja es un concepto distinto y de naturaleza más operativa. Se calcula mensualmente y combina ingresos y pagos del ejercicio en curso con los de presupuestos de ejercicios ya cerrados que se materializan efectivamente en tesorería durante el año. La divergencia más clara entre estos dos conceptos queda patente en el caso de los intereses por deuda, cuyo tratamiento temporal es distinto: el criterio de caja los imputa presupuestariamente en el momento del pago efectivo, mientras que la contabilidad nacional los prorratea a lo largo de toda la vida del título de deuda.\footnote{
Sin embargo, la casuística es muy grande. Por ejemplo, si un tribunal anula un impuesto o una tasa y obliga a la devolución de lo ingresado. El criterio de caja contabilizará los pagos materiales que se hacen por ello, determinando si cabe la necesidad de endeudamiento. Sin embargo, para introducirlo en contabilidad nacional, el primer paso exige que Eurostat decida cómo va a realizar, por ejemplo si se modifica el resultado de los años en que se cobró o si todo se imputa al año en que se dictó la sentencia o a los años en que el Estado la ejecuta. 

En todo caso, las normas contables no son objeto de este tema.
}

Como podemos intuir, las necesidades de financiación del Estado, derivadas de este déficit de caja no financiero y de la variación neta de activos financieros, son las que se cubren mediante emisión de deuda. Así, el saldo presupuestario y, con más frecuencia, el saldo de caja son la magnitud que determina, mes a mes, cuánta deuda necesita emitir el Tesoro Público para hacer frente a sus compromisos de pago inmediatos, con independencia de a qué ejercicio presupuestario correspondan en origen. ¿Cómo? Literalmente, computando, mensual y anualmente, los ingresos y pagos presupuestados en el año en curso y anteriores, si estos se materializan ahora. El resultante determina las emisiones de deuda a lo largo del año.

La emisión neta prevista para 2023 era de 70.000 millones de euros, pero la emisión bruta, que incluye la refinanciación de la deuda que vence ese mismo año, ascendía a 256.930 millones. Esta diferencia tan sustancial entre emisión neta y bruta es, en sí misma, un dato revelador de la magnitud de la gestión de tesorería que el Tesoro Público debe sostener año tras año, muy por encima de lo que la simple cifra del déficit anual del ejercicio en curso podría sugerir a primera vista. Por tanto, gran parte de la actividad de emisión de deuda española no responde, en absoluto, a necesidades nuevas de financiación, sino a la pura mecánica de sustituir, sistemáticamente, deuda antigua que vence por deuda nueva que la sustituye.

\subsection{Organismos Autónomos}
Los últimos PGE (2023) incluían $57$ Organismos Autónomos, con un presupuesto conjunto de $44,7$mm (3,2\% PIB). Como sabemos ([\authorfont{Ver Tema 4.B.21}]), la existencia de estos organismos responde a una necesidad de especialización, permitiendo la LRJSP la personificación instrumental para dotar de gestión especializada a una actividad de naturaleza administrativa sin salir del control jerárquico ministerial. 

En lo que respecta al gasto de estos organismos, sin duda el capítulo más importante es el de Transferencias Corrientes (IV), alcanzando un $81,2\%$ del total. Muy por debajo quedaría el segundo, Transferencias de Capital (VII), con un $4,5\%$. Desde otra perspectiva, más orgánica, los Organismos más importantes son el Servicio Público de Empleo Estatal (SEPE), el Fondo Español de Garantía Agraria (FEGA), la Mutualidad General de Funcionarios Civiles del Estado (MUFACE), La Jefatura Central de Tráfico, el Instituto Social de las Fuerzas Armadas (ISFAS), el Fondo de Garantía Salarial (FOGASA) y las Confederaciones Hidrográficas. No obstante, hay que tener presente que SEPE y FEGA abarcan el $79\%$ del gasto total.

En lo relativo al ingreso, ascendió a un total de $45.9$mm de euros ($2,4\%$ PIB). La principal fuente de ingresos fue el capítulo de Impuestos directos y cotizaciones sociales (Capítulo I), con el $66,5\%$ del total. Como podemos suponer, fundamentalmente por las cotizaciones que ingresan el SEPE, MUFACE, ISFAS, FOGASA y MUJEJU, cada uno en su ámbito subjetivo/material de aplicación. Los recursos propios de los Organismos Autónomos  alcanzan el $71,5\%$ de sus ingresos totales, y el Capítulo I se ve complementado para ello por el Capítulo III, Tasas, y los de carácter patrimonial (Capitulo V). Sin embargo, el segundo capítulo por importancia es el de Transferencias Corrientes (IV), proporcionando el $19.,6\%$ de los recursos totalel. Debido a las transferencias que reciben, esencialmente el SEPE, por pare de la Administración del Estado.\footnote{%
    En el caso del SEPE, el Estado le transfiere fondos para pagar las prestaciones no contributivas de desempleo. Al igual que en el caso de los gastos, esta partida puede fluctuar mucho a lo largo de los años, según la situación económica y según que el SEPE disponga o no de remanentes de ejercicios anteriores.
}

Por último, destacan los ingresos por Transferencias de capital, cerca del $65\%$ del total de ingresos. El principal receptor es el FEGA, que los recibe del Fondo Europeo Agrícola de Desarrollo Rural (FEADER) para medidas de desarrollo rural bajo el paraguas de la PAC. También destacan por recibir este tipo de transferencias el Instituto de Salud Carlos III y las Confederaciones Hidrográficas.

Veamos más en especial los dos Organismos que concentran el $79\%$ del gasto total del subsector: SEPE y FEGA.

\subsubsection{Servicio Público de Empleo Estatal (SEPE): Clasificación presupuestaria}
El SEPE es, con el $61,9\%$ del gasto de todos los Organismos Autónomos ($27,8$mm, el $2\%$ del PIB), la entidad de mayor peso del subsector. Sus gastos, en el PGE23, se concentran casi por completo ($98,3\%$) en transferencias corrientes.El $78\%$ de estas transferencias corresponden a prestaciones de desempleo, $21,3$mm. Además, las prestaciones se subdividen en: prestaciones contributivas (13mm), subsidios por desempleo (6,6mm), subsidios del SEASS, renta activa de reinserción (0,7mm) y gastos de funcionamiento (0,4mm). 

Por otro lado, la partida de mayor gasto son las políticas activas de empleo (6,2mm). Son el conjunto de servicios y programas de orientación, empleo y formación para el empleo en el ámbito laboral dirigidas a mejorar las posibilidades de acceso al empleo, por cuenta ajena o propia, de las personas desempleadas, al mantenimiento del empleo y a la promoción profesional de las personas ocupadas y al fomento del espíritu empresarial y de la economía social. Dentro de las políticas activas, destacan: 
\begin{lnum}
    \item \textbf{Inserción} (4,1mm). El programa fomento de la inserción y estabilidad laboral. que tiene por objetivo incrementar el empleo promoviendo la integración de los desempleados en el mercado laboral y la permanencia en el mismo de los trabajadores, especialmente de colecti vos más desfavorecidos.
    \item \textbf{Formación} (2,1mm). La formación profesional para el empleo, que engloba una serie de medidas que tienen por objetivo impulsar entTe los trabajadores ocupados y los desempleados una form ación profes ional que contribuya a mejorar su cualificación y la competitividad de las empresas.
    \item \textbf{Prestación} (6,14mm). El programa prestaciones económicas por cese de actividad, que tiene por objeto regular un sistema de protección por cese de actividad de los trabajadores autónomos, así como impulsar medidas de fo rmac ión, orientación y promoción de la actividad emprendedora dotándose al efecto con 2 millones de euros.
\end{lnum}

Conviene volver a llamar la atención sobre las diferencias entre los datos de liquidación presupuestaria y los de contabilidad nacional. En este último caso, los datos de la clasificación COFOG no coinciden con los de liquidación presupuestaria. Uno de los motivos es que la COFOG incluye las prestaciones por cese de actividad de la Seguridad Social, mientras que en contabilidad presupuestaria se consignan los del SEPE (y Seguridad Social, pero por otro lado).

Por otro lado, es necesario contextualizar las cifras del SEPE, por su naturaleza, con su enorme volatilidad histórica. En el peor momento de la crisis financiera de 2008-2012, el gasto en prestaciones alcanzó 31,5mm en 2012. Posteriormente, con la irrupción de la Covid-19 en 2020, se disparó hasta 34,3mm, el máximo histórico, arrastrado por el mecanismo de los ERTE (Expedientes de Regulación Temporal de Empleo), que canalizaron a través del SEPE un volumen de prestaciones sin precedente en la historia de la institución.\footnote{%
    España, con el ERTE, optó por un instrumento de naturaleza muy similar al Coronavirus Job Retention Scheme británico o al Kurzarbeit alemán: mantener el vínculo laboral formal entre empresa y trabajador, financiando públicamente una parte sustancial del salario durante la suspensión temporal de la actividad, en lugar de dejar que la relación laboral se extinguiera y el trabajador pasara a depender íntegramente de la prestación ordinaria por desempleo. Una decisión de diseño que explica por qué, pese al volumen extraordinario de gasto de 2020, la tasa de desempleo estructural española no llegó a dispararse con la misma intensidad que en crisis anteriores.
}

\subsubsection{Fondo Español de Garantía Agraria (FEGA)}
El Fondo Español de Garantía Agraria (FEGA), segundo organismo por volumen de gasto, no es una institución de diseño puramente español.

Como decíamos, es el organismo pagador a través del cual España gestiona y canaliza los fondos de la Política Agrícola Común (PAC). La PAC se articula en torno a dos fondos: el FEAGA (Fondo Europeo Agrícola de Garantía), que financia fundamentalmente los pagos directos a los agricultores y las medidas de mercado, y el FEADER (Fondo Europeo Agrícola de Desarrollo Rural), que cofinancia, junto con el propio Estado miembro, programas de desarrollo rural. Las transferencias corrientes del FEGA ($5,8$mm) corresponden a las subvenciones a la producción agraria cofinanciadas por el FEAGA, mientras que sus transferencias de capital (1,6mm) corresponden al desarrollo rural cofinanciado por el FEADER.

La PAC nació, en los años sesenta, con un diseño que generó, ya en los años ochenta, los tristemente célebres excedentes (las montañas de mantequilla y los lagos de vino que la propia prensa europea de la época popularizó) por incentivar una sobreproducción estructural desconectada de la demanda real del mercado. La reforma MACSHARRY de 1992 inició la transición desde ese modelo de sostenimiento de precios hacia el actual sistema de pagos directos desacoplados de la producción, complementado desde entonces con sucesivas rondas de reforma que han ido incorporando condicionalidad medioambiental, exigiendo a los agricultores beneficiarios el cumplimiento de determinadas prácticas de sostenibilidad ambiental. 

Las cifras del FEGA que aparecen en los PGE españoles son, en este sentido, la manifestación contable de una de las políticas más debatidas de toda la historia de la integración europea, un debate que trasciende, con mucho, el ámbito puramente presupuestario.

### El resto del subsector: tres instituciones que revelan, cada una, una lógica propia

FOGASA, el Fondo de Garantía Salarial, responde a una lógica completamente distinta. Garantiza a los trabajadores el cobro de los salarios pendientes cuando su empresa resulta insolvente, en desarrollo de una exigencia mínima que el propio Derecho de la Unión Europea impone a todos sus Estados miembros, de modo que su existencia responde a una obligación de armonización social europea y no a una decisión puramente doméstica. 

Y las Confederaciones Hidrográficas responden a una tercera lógica, de naturaleza casi territorial. Los principales ríos españoles atraviesan varias Comunidades Autónomas a la vez, de modo que su gestión no puede quedar, sin generar conflictos de coordinación evidentes, en manos de una única Administración autonómica, justificando la existencia de un organismo de ámbito supraautonómico, dependiente del Estado, específicamente diseñado para gestionar un recurso cuya propia geografía física desborda cualquier frontera autonómica.

\subsection{Seguridad Social}
Antes de entrar en las cifras del PGE23, conviene fijar de dónde procede la propia institución, porque su historia explica buena parte de su complejidad organizativa actual. 

Dada su importancia económica, la Seguridad Social tiene asignada en la clasificación orgánica de los PGE una Sección, la Sección 60, distinta de la asignada al Ministerio de Inclusión. Migraciones y Seguridad Social. en el que está encuadrada.

En 2023, los gastos totales presupuestados en la sección de Seguridad Social ascendieron a $204,2$mm, el $14,7\%$ del PIB. Por el lado de ingresos, que suma $204,2$mm ($14,7\%$ del PIB).\footnote{%
    Siguiendo la clasificación económica para gastos, el único capítulo verdaderamente reseñable es el de Transferencias Corrientes, que representa el $95,2\%$ del total. Por lo tanto, la clasificación económica tiene escaso interés.

Según la clasificación económica para ingresos, los únicos capítulos relevantes son las Cotizaciones Sociales (Capítulo 1), con el $74,5\%$, las Transferencias Corrientes, con el $19\%$ y los Pasivos Financieros, con el $4,9\%$.
} 

\subsubsection{Clasificación orgánica}
El ámbito institucional del presupuesto de la Seguridad Social abarca un conjunto de entidades que pueden clasificarse en tres grupos diferenciados: entidades gestoras, servicios comunes y mutuas colaboradoras con la Seguridad Social y sus centros mancomunados. 

Las \textbf{entidades gestores} son cuatro. En primer lugar, el Instituto Nacional de la Seguridad Social (INSS), que es la entidad a la que se le encomienda la gestión y administración de las prestaciones económicas del sistema de la Seguridad Social, con excepción de las que corresponden al IMSERSO y al ISM. En segundo lugar, el Instituto Nacional de Gestión Sanitaria (INGESA), heredero del Instituto Nacional de la Salud (INSALUD), que gestiona las prestaciones sanitarias en el ámbito de las ciudades autónomas y las necesidades del Centro Nacional de Dosimetría de Valencia.\footnote{%
    Esta función tan territorial tiene una razón: las competencias en materia de Sanidad han sido transferidas a las Comunidades Autónomas, a excepción de las que siguen estando bajo la gestión del INGESA.
} En tercer lugar tendríamos el Instituto de Mayores y Servicios Sociales (IMSERSO), adscrito al Ministerio de Derechos Sociales y Agenda 2030, que gestiona las pensiones no contributivas de invalidez y jubilación, así como de los servicios complementarios de las prestaciones de la Seguridad Social y la gestión de planes, programas y servicios de ámbito estatal para personas mayores y para personas en situación de dependencia. En cuarto, y último lugar, el Instituto Social de la Marina, que es el organismo específico y unitario de gestión de la protección de los trabajadores del mar. Le corresponde la gestión, administración y reconocimiento del derecho a las prestaciones económicas, asistencia sanitaria y servicios sociales del régimen especial de la Seguridad Social de este colectivo. 

Los \textbf{servicios comunes} son entidades, con o sin personalidad jurídica propia, que desarrollan tareas comunes entre todos los entes que integran el sistema de la Seguridad Social. Son la Tesorería General de la Seguridad Social, la Gerencia de Información de la Seguridad Social y el Servicio Jurídico de la Administración de la Seguridad Social. 

Por último, las \textbf{mutuas colaboradoras} con la Seguridad Social son asociaciones privadas de empresarios, sin ánimo de lucro, que colaboran en la gestión de las tareas encomendadas al Sistema de Provisión y Seguridad Social. Por ejemplo, la gestión de prestaciones económicas y de asistencia sanitaria por accidente de trabajo y enfermedad profesional, la gestión de las prestaciones económicas por incapacidad temporal, por riesgo en el embarazo, por cuidado de menores etc.

En esta clasificación orgánica del Presupuesto de la Seguridad Social, el INSS gestiona el $88,40\%$, seguido de las Mutuas con el $8,5\%$, y de forma residual el reparto entre las tres restantes, INGESA e IMSERSO. Por lo que respecta al ingreso, \textcolor{red}{\textbf{CUMPLIMENTAR}...}





\subsubsection{Clasificación Funcional}
Conforme a la clasificación funcional, los gastos se dividen en cuatro areas principales: Prestaciones Económicas, Asistencia Sanitaria, Servicios Sociales y Servicios Comunes. Cada área es gestionada en mayor o en menor medida por los cuatro centros gestores, salvo los Servicios Comunes.

Como la mayor parte del presupuesto de la Seguridad social son transferencias corrientes, el área de gasto más importante, con diferencia, son la de las Prestaciones Económicas, que ocupan el $93,9\%$ del presupuesto. A continuación, estarían los Servicios Sociales, con el $1,86\%$ y la Asistencia Sanitaria, con el $0,89\%$. Por su importancia, nos centraremos en las prestaciones económicas.

Dentro de esta categorías, vemos transferencias por garantías de diferencia índole. Las pensiones contributivas suman el $85,8\%$ y pueden ocasionarse por motivos de jubilación, invalidez, viudedad, orfandad y en favor de familiares. Además, incluyen complementos de mínimos que son financiados directamente por la Administración Central (Estado). Por otro lado, las pensiones no contributivas, que figuran en el presupuesto del IMSERSO, suponen el $1,4\%$ de las prestaciones económicas y son prestaciones de invalidez y de jubilación (de índole económica, médica, farmacéutica y de servicios sociales) para personas que se encuentren en situación de necesidad al carecer de rentas o ingresos. El número de personas beneficiarias asciende a las $428$ mil, y se financian mediante aportaciones directas de la Administración Central. La prestación de Incapacidad Temporal integra los subsidios para compensar las consecuencias económicas derivadas de la situación de baja laboral por enfermedad o accidente, son el $6,1\%$ de las prestaciones económicas. Las prestaciones por nacimiento y cuidado del menor y otras están vinculadas al cese temporal en el trabajo por nacimiento y cuidado del menor, riesgo durante el embarazo y la lactancia natural, cuidado de menores afectados por cáncer u otra enfermedad grave y corresponsabílidad en cuidado del lactante, suponen el $1,9\%$ del total. Las prestaciones de Atención a la Dependencia están integradas en el Sistema para la Autonomía y Atención a la Dependencia, suponen el $1,8\%$ del total y se complementan con prestaciones de asistencia sanitaria y residencial. Por último, el Ingreso Mínimo Vital, que se configura como el derecho subjetivo a una prestación económica que garantiza un nivel mínimo de renta a quienes se encuentren en situación de vulnerabilidad económica con el objetivo de garantizar una mejora de oportunidades reales de inclusión social y laboral de las personas beneficiarias, supone el $1,4\%$ del total.


Desde la óptica de ingresos, es la clasificación económica la que nos proporciona más información relavante. 

La aportación más cuantiosa corresponde a las Cotizaciones Sociales, donde el $94,2\%$ procede de empresas y trabajadores. El resto, las cotizaciones de desempleados y beneficiarios por cese de actividad, son residuales.\footnote{%
    Desde una óptca funcional, el régimen que aporta más cotizaciones es el Régimen General de la Seguridad Social, con el $76,5\%$. Después, el Régimen especial de Autónomos, con el $7,9\%$, las cotizaciones por accidentes de trabajo y enfermedades profesionales, con el $7,6\%$, y las transferencias del SEPE y cotizaciones de desempleados con el $5,5\%$ en conjunto.
} Las Transferencias corrientes, por su parte, proceden del Estado y ascienden a $38,8$mm ($2,8\%$ del PlB). Una parte sustancial de las transferencias se dedica a financiar prestaciones no contributivas, las cuales se califican como \textit{gastos impropios} de la Seguridad Social. El resto, por importe de $19,9$mm, se realiza en cumplimiento de la recomendación primera del Pacto de Toledo de 2020, con el objeto de garantizar la sostenibilidad del sistema a medio y largo plazo.\footnote{%
    Con esta cantidad se financian, por ejemplo: las prestaciones por nacimiento y cuidado del menor, bajas de maternidad y paternidad, reducciones en la cotización del Régimen Especial del Mar, la integración de lagunas de cotización, los complementos de pensiones contributivas para progenitores de dos o más hijos, etc.
} No obstante, las razones por las cuales el Estado transfiere fondos es una larga lista de conceptos, como la asistencia sanitaria del INGESA, la financiación del mínimo garantizado a través del IMSERSO, la cobertura de los complementos para pensiones mínimas, la financiación del Ingreso Mínimo Vital, entre otras.

Por último, los ingresos por Pasivos Financieros reflejan los préstamos, sin interés, que el Estado concede a la Seguridad Social con el fin de garantizar el equilibrio presupuestario. La razón de contabilizarlos como préstamos es, de nuevo, que no implican la traslación del déficit de una a otra Administración. Como sabemos, la transferencia implicaría un aumento del déficit del Estado, imposible de realizar por los compromisos fiscales. Por lo tanto, instrumentarlos como préstamos implica no trasladar el déficit en términos de contabilidad presupuestaria. 

No obstante, este recurso analítico no impide que sí se contabilicen como déficit al considerar las Administraciones Públicas en su conjunto. Por ejemplo, en contabilidad nacional.

Tratándose de la Administración Central, el problema es el gasto en pensiones. No obstante, a pesar de que el sistema sanitario y el de cuidados de dependencia están transferidos. el Estado, como garante de unas prestaciones suficientes en todo el territorio nacional. también realiza transferencias anuales (corrientes y de capital) a las CCAA para cubrir las necesidades de financiación que éstas afrontan en estas áreas, que no estén cubiertas suficientemente por el sistema de financiación autonómico.


\subsection{Entidades con presupuesto limitativo}
Antes de abordar el presupuesto consolidado, nos refirmos someramente al resto de entidades del sector público administrativo con presupuesto limitativo, que se integrarán en la versión consolidada.\footnote{%
    \textbf{NOTA.-} \textcolor{red}{Aquí hay que incluir la justificación de por qué no vamos a tratar las entidades con presupuesto estimativo, que se han trasladado al Anexo}. [\authorfont{Ver Anexo \ref{anx:presupuesto_estimativo}}]
}

Esta compuesto por un conjunto de veinte entidades, con un gasto conjunto de $9,5$ millones de euros, del que destacan la CNMC ($3,7$ millones), la $AEI$ (1,4 millones), la $AEAT$ (1,4 millones) y el $CSIC$ (1,2 millones). 

Conviene verificar con cuidado las cifras contra la fuente primaria.

Los más importantes son los de personal. con 2. 1 millones • Las transferencias corrientes corresponden en su mayor paiie a la· AECJD41, por importe de 550 millones, principalmente para ayudas al desarrollo
• En el capítulo de inversiones destaca el CSlC. para infraestructuras de investigación
• En cuanto a h·ansferencias ele capital. destacan la CNMC. siendo la principal actuación financiar los costes del sector eléctrico; y la AEI. fundamentalmente para la participación en el Fondo para la Investigación Científica y el Desarrollo Tecnológico, y para el CSIC.

Los ingresos de estas entidades equivalen a sus gastos (9,5 milloues) y proceden sobre todo de transferencias corrientes (3 millones) y de capital (5,7 millones) del Estado42. • En ingresos por transferencias corrientes destacan la AEAT, la AECID y el CSIC • En transferencias de capital destaca la CNMC

\subsubsection{CNMC: el "macrorregulador" que solo España ha llevado tan lejos}
La Comisión Nacional de los Mercados y la Competencia, creada por la Ley 3/2013, es la entidad de mayor peso de todo este grupo.

Su creación supuso la fusión de seis organismos reguladores sectoriales hasta entonces independientes entre sí. La Comisión Nacional de la Competencia (CNC), la Comisión Nacional de Energía (CNE, 1995), la Comisión del Mercado de las Telecomunicaciones (CMT, 1996), la Comisión Nacional del Sector Postal, el Consejo Estatal de Medios Audiovisuales, y el Comité de Regulación Ferroviaria y la Comisión de Regulación Económica Aeroportuaria; estos dos últimos, fusionados a su vez poco antes en un único órgano. 

Lo verdaderamente singular de esta operación es que España es el único Estado miembro de la Unión Europea que ha llevado la unificación regulatoria tan lejos. Los Países Bajos, el único otro caso de fusión regulatoria, se limitó a integrar la autoridad de competencia con los reguladores de telecomunicaciones y consumo, sin llegar a incorporar energía, transporte ferroviario y aéreo, sector postal y medios audiovisuales bajo un único paraguas institucional como sí hizo España.

Esta ambición integradora no estuvo exenta de una controversia jurídica real y bien documentada. Los presidentes cesados de la CNE y la CMT, al ver terminado anticipadamente su mandato con motivo de la fusión, recurrieron ante el Tribunal Supremo, que planteó, en auto de 3 de julio de 2015, tres cuestiones prejudiciales ante el TJUE sobre la compatibilidad de la propia CNMC con el Derecho comunitario, que exigía a los Estados miembros garantizar reguladores sectoriales especializados y dotados de independencia y estabilidad de mandato propias para sectores como las telecomunicaciones. 

La pregunta de fondo que el litigio planteaba es si puede un macrorregulador único, que fusiona funciones y remueve a sus responsables sectoriales previos en el proceso, seguir ofreciendo el mismo grado de independencia especializada que el Derecho europeo exige para cada sector regulado individualmente.

Las transferencias de capital que la CNMC gestiona dentro de los PGE corresponden, en gran medida, a la cobertura de los costes regulados del sistema eléctrico y gasista, un terreno que conecta con el histórico problema español del déficit de tarifa eléctrica: durante años, el precio regulado de la electricidad en España se fijó por debajo de los costes reales, generando un déficit acumulado que hubo de titulizarse y financiarse a largo plazo, y cuya gestión y supervisión corresponde a la CNMC como heredera de las funciones de la antigua CNE.

\subsubsection{AEI: el superviviente inesperado del modelo de agencias que el resto del sistema abandonó}
La Agencia Estatal de Investigación, creada en 2015 al amparo de la Ley 28/2006 de Agencias Estatales, merece una mención específica por una razón que conecta directamente con el debate que ya cerramos en el Tema 21: recordarás que la propia Ley 40/2015 eliminó, casi en el mismo momento histórico, la figura de la agencia estatal del catálogo general del sector público institucional, por su incompatibilidad práctica con la disciplina fiscal reforzada tras la crisis.

La AEI es, sin embargo, uno de los pocos ejemplos que ha sobrevivido como agencia estatal plenamente operativa bajo ese régimen jurídico, sugiriendo que el modelo de gestión por agencia —con la autonomía operativa y la financiación plurianual que llevaba aparejadas— encuentra un encaje funcional más natural en la gestión de la investigación científica, donde la propia naturaleza de los proyectos financiados exige horizontes de planificación más largos y menos sujetos al ciclo presupuestario anual estricto, que en otras funciones administrativas donde el experimento de las agencias, como ya vimos con el caso del propio SEPE, resultó menos sostenible.

### La AEAT: la agencia tributaria como pionera del modelo de gestión por resultados en España

[Probable] La Agencia Estatal de Administración Tributaria, creada por la Ley 31/1990, de Presupuestos Generales del Estado para 1991 —una de las primeras "agencias" del sistema español, anterior incluso a la propia Ley de Agencias de 2006—, responde a la lógica de dotar a la función recaudatoria, de complejidad técnica creciente y de importancia estratégica evidente para el conjunto del sistema fiscal español que hemos recorrido a lo largo de toda esta serie de temas, de una gestión más ágil y más especializada que la que permitiría una Dirección General ordinaria integrada en la estructura jerárquica convencional del Ministerio de Hacienda.

### El CSIC: de la Junta para Ampliación de Estudios a la Ciencia bajo tutela del nuevo régimen

[Seguro] El Consejo Superior de Investigaciones Científicas, principal organismo público de investigación español, fue creado en 1939, en los primeros meses tras el final de la Guerra Civil, y su fecha de nacimiento no es un dato anecdótico: sustituyó de forma explícita a la Junta para Ampliación de Estudios e Investigaciones Científicas (JAE), creada en 1907, la institución que había vertebrado, durante las tres décadas anteriores a la Guerra Civil, buena parte de lo que la historiografía cultural española conoce como la Edad de Plata de la ciencia y la cultura españolas —becas de estudio en el extranjero, residencias de estudiantes, laboratorios de investigación de primer nivel internacional para su época—. [Probable] La sustitución de la JAE por el CSIC en 1939 no fue una simple continuidad institucional bajo nuevo nombre: fue una ruptura deliberada, promovida explícitamente por el nuevo régimen franquista como instrumento de recatolización y de control ideológico de la actividad científica española, bajo la tutela directa del propio Ministerio de Educación Nacional de la época y con una impronta confesional muy marcada en sus primeros años de funcionamiento —un episodio que conviene conocer porque revela, una vez más, algo que ya hemos visto repetirse a lo largo de todo este proyecto: las instituciones del sector público español que hoy consideramos técnicas y neutrales rara vez lo fueron en su origen histórico, y entender esa genealogía ayuda a comprender mejor su configuración actual—. El CSIC de hoy, evidentemente, es una institución científica plenamente moderna y alineada con los estándares internacionales de investigación, pero su fecha de nacimiento sigue siendo, ochenta y cinco años después, testimonio de una ruptura histórica muy concreta dentro de la ciencia española del siglo XX.

### La AECID y el compromiso internacional del 0,7%

[Probable] Cierro con la Agencia Española de Cooperación Internacional para el Desarrollo, principal receptora de las transferencias corrientes de este grupo de entidades, gestora de la ayuda oficial al desarrollo española conforme a la Ley 23/1998, de Cooperación Internacional para el Desarrollo. Conviene situar su actividad dentro del marco de referencia internacional que le da sentido: el compromiso, asumido por Naciones Unidas desde los años setenta, de que los países desarrollados destinen el 0,7% de su Renta Nacional Bruta a ayuda oficial al desarrollo —un objetivo que España, como la inmensa mayoría de los países donantes desarrollados, ha incumplido de forma sistemática durante décadas, situándose su ayuda oficial al desarrollo efectiva muy por debajo de ese umbral de referencia internacional en la práctica totalidad de los ejercicios presupuestarios recientes—, un dato que conviene conocer como contexto de fondo de las cifras, modestas en términos relativos del presupuesto total del Estado, que la AECID gestiona año tras año.


\clearpage
\section{Contabilidad presupuestaria: Visión conolidada}
\puntosclave{
    Los PGE consolidados permiten conocer la acción de la Adm i nistración Central en su conjunto. 

    El gasto total asciende a $553,5$mm, casi el $40\%$ del PIB. Con el PRTR es un $5,4\%$ más.

    El gasto de operaciones financieras es de $425,6$mm ($30,6\%$ del PIB). Con el PRTR es un 6 por ciento más.

    Por subsectores, por orden de importancia, son: el Estado ($47,3\%$), la Seguridad Social ($41,3\%$), OOAA ($9,4\%$) y Resto de entidades ($2\%$).

    En los ingresos, los más importantes son las cotizaciones sociales ($44,6\%$), los impuestos directos ($20\%$) y los impuestos indirectos ($15,4\%$).
    
    Hay $27$ políticas de gasto, agrupadas en cinco grandes áreas de gasto: las Actuaciones de protección y promoción social son el $50.5\%$ y si se le suma la Producción de bienes públicos de carácter preferente, se obtiene un gasto social del $55.7\%$. Siguen en importancia las Actuaciones de carácter general con el $31\%$, debido a las transferencias a otras AAPP y las Actuacionés de carácter económico, que suponen un $8\%$ del gasto.
}

Si hemos visto, entidad por entidad, que el Estado transfiere masivamente a la Seguridad Social, a los Organismos Autónomos y al resto de entidades limitativas, y que estas, a su vez, reciben esas transferencias como ingreso propio, entonces sumar sin más los presupuestos de cada una nos daría una cifra inflada y engañosa del tamaño real del sector público administrativo estatal. 

Debemos, por tato, cerrar presentando la organización presupuestaria que da sentido en retrospectiva a todo lo presentado hasta ahora: la visión consolidada. Se trata, en esencia, de la respuesta técnica al problema de agregación y, antes de aplicarla con las cifras observadas, conviene entender por qué hace falta, qué gana y qué pierde el análisis al construirla, y cómo se construye exactamente.

\subsection{Porqué y el cómo de la consolidación presupuestaria}
La literatura técnica internacional de estadísticas de finanzas públicas (Manual del del FMI, GFSM) distingue tres operaciones que el lenguaje coloquial confunde con frecuencia bajo la etiqueta genérica de sumar cuentas: la agregación, la consolidación y el neteo. 

La \textbf{agregación} es la simple suma de las cifras brutas de varias unidades, preservando cualquier relación de crédito o débito que exista entre ellas. Por ejemplo, si sumamos de forma agregada el presupuesto del Estado y el de la Seguridad Social, cualquier transferencia entre ambos aparece registrada dos veces, una como gasto de quien la realiza y otra como ingreso de quien la recibe, sin que la suma elimine esa duplicidad. La \textbf{consolidación}, en cambio, es la operación que sí elimina esas posiciones y esos flujos recíprocos entre las unidades agrupadas, antes de sumar el resto, exactamente lo que necesitamos para que el gasto total del sector público estatal no cuente dos veces el mismo euro. El \textbf{neteo}, la tercera operación, es distinta de ambas. Consiste en presentar solo el saldo neto de determinadas partidas de signo contrario dentro de una misma unidad, sin llegar a ser propiamente ni agregación ni consolidación. 

Esta distinción tiene consecuencias prácticas reales. Cualquier cifra de gasto público total que encontremos en un informe o en un medio de comunicación puede estar construida bajo cualquiera de estas tres lógicas, y solo la consolidación ofrece la magnitud que realmente interesa para conocer cuánto extrae o inyecta el sector público, en su conjunto, en el resto de la economía.

### La genealogía: de nuevo, Richard Stone, y una diferencia importante que conviene no pasar por alto

La técnica de consolidación tiene la misma raíz intelectual que el Sistema de Cuentas Nacionales (STONE), pero con un matiz que revela una diferencia real entre los dos grandes sistemas estadísticos internacionales.

El Sistema de Cuentas Nacionales (SNA/SEC) no aplica, por regla general, la consolidación. Si una unidad de gobierno posee deuda pública emitida por otra unidad de gobierno del mismo sector, ambas posiciones (el activo financiero de quien la posee y el pasivo de quien la emitió) se mantienen registradas por separado en el balance del sector. El sistema de Estadísticas de Finanzas Públicas (GFS), en cambio, sí aplica consolidación plena, precisamente con el objetivo declarado de mostrar las actividades del sector de las administraciones públicas, o de cualquier otro agrupamiento de unidades, como si existiera una única unidad. 

Esta diferencia explica por qué debeos dedicar un apartado entero a la visión consolidada en términos de contabilidad presupuestaria, y tratar de forma separada la contabilidad nacional ([\authorfont{Ver Tema 4.B.23}]). No son, simplemente, dos niveles de agregación territorial distintos, son, en algunos aspectos técnicos de fondo, dos convenciones contables que no consolidan exactamente de la misma manera.

Consolidación intrasectorial frente a intersectorial: por qué este tema es, literalmente, el paso previo al siguiente

El propio GFSM 2014 distingue dos niveles de consolidación. La consolidación intrasectorial elimina los flujos y posiciones que se producen dentro de un mismo subsector o unidad institucional. Por ejemplo, dentro de la administración central de un país. La consolidación intersectorial elimina, en cambio, los flujos entre distintos subsectores o sectores. Por ejemplo, entre la administración central, las administraciones regionales y las administraciones locales. 

Y, además, el propio manual recomienda expresamente que la consolidación intrasectorial se realice antes que la intersectorial. Esto significa que la visión consolidada de los PGE que vamos a construir (consolidando Estado, Organismos Autónomos, Seguridad Social y resto de entidades limitativas, todas ellas dentro del mismo subsector de Administración Central) es el ejercicio de consolidación que la normativa identifica como el paso técnico previo e indispensable antes de abordar la consolidación intersectorial completa (Administración Central, Comunidades Autónomas, Entidades Locales y Seguridad Social como subsectores distintos) que constituye, precisamente, el objeto central del siguiente tema ([\authorfont{Ver Tema 4.B.23}]). 

### Por qué la comunidad estadística internacional exige esta operación: el argumento de la comparabilidad

El trabajo de O'CONNOR, WEISMAN y WICKENS (2004), documento técnico del FMI sobre la consolidación del sector de las administraciones públicas, formula con precisión el argumento de fondo que justifica la consolidación, y no la mera agregación, como el estándar exigible para cualquier presentación seria de cifras de gasto público: la consolidación elimina los efectos sobre los agregados de las diferencias en los arreglos administrativos entre países. 

Detengamomos para ver qué significa esto. Imaginemos dos países, que organizan la misma función pública, la gestión del empleo, de formas institucionalmente distintas. Uno puede gestionarla directamente desde un departamento ministerial sin personalidad jurídica propia, mientras que otro la gestiona a través de un organismo autónomo separado que recibe transferencias del Estado central. Si ambos países presentaran sus cifras de gasto público sin consolidar, el segundo país mostraría, de forma puramente artificial, una cifra de gasto total mayor que el primero, aunque ambos estuvieran gastando exactamente lo mismo en la misma función. ¿Por qué? Evidentemente, porque contaría dos veces la transferencia interna que el primer país nunca necesitó registrar, al no existir esa separación institucional. 

La consolidación es, en este sentido, el instrumento que permite comparar el tamaño real del sector público entre países con arquitecturas administrativas internas radicalmente distintas, neutralizando lo que de otro modo sería puro artefacto organizativo y no diferencia económica real.

### Qué problemas resuelve, con precisión, y por quién se aplica

Por tanto, resuelve la doble contabilización de las transferencias internas entre las entidades del perímetro consolidado, ofreciendo así la magnitud real, comparable internacionalmente, y comparable con el PIB, de la actividad del sector público frente al resto de la economía. En España, esta consolidación la realiza la propia Dirección General de Presupuestos, dentro de la documentación oficial del Proyecto de Ley de PGE (\textit{serie amarilla}), bajo la supervisión última de la IGAE en su función de elaboración de estadísticas y de rendición de cuentas.

### Qué problemas NO resuelve: dos límites que conviene conocer con la misma honestidad que hemos aplicado en cada bloque de este proyecto

No obstante, conviene también advertir de dos límites que presenta la visión consolidad. 

El primer límite, como ya anticipamos, es que la consolidación no captura el peso económico real de las entidades con presupuesto estimativo, que quedan fuera del perímetro consolidado que estamos construyendo.\footnote{Por ejemplo, en España, ENAIRE, ADIF, RENFE-Operadora, ICO, SEPI.
} Esto genera una vulnerabilidad conceptual, ya que si un Gobierno quisiera, por la razón que fuera, que la cifra consolidada de los PGE pareciera más contenida de lo que realmente es, podría, en teoría, canalizar una proporción creciente de actividad económica pública a través de entidades con presupuesto estimativo sin que ese desplazamiento se refleje jamás en la cifra consolidada que constituye el titular mediático y político de cada Ley de Presupuestos.\footnote{%
    Por ejemplo, aumentando su capital, encomendándoles nuevas funciones, financiándolas mediante aportaciones patrimoniales del Capítulo VIII en lugar de mediante gasto directo;
} Es la misma lógica que aplica a los fondos extrapresupuestarios como técnica de opacidad ([\authorfont{Ver Tema 4.B.21}]). La consolidación resuelve el problema de la doble contabilización dentro de un perímetro, pero no resuelve la pregunta previa y más importante de si ese perímetro está bien trazado.

El segundo límite emerge del tratamiento de las modificaciones presupuestarias. La cifra consolidada que vamos a presentar es la del presupuesto inicialmente aprobado, no la del presupuesto finalmente ejecutado tras las sucesivas transferencias, ampliaciones, generaciones y créditos extraordinarios ([\authorfont{Ver Tema 4.B.21}]). La consolidación, por tanto, no puede resolver el problema de la divergencia temporal entre lo que las Cortes aprueban en diciembre y lo que el sector público estatal termina gastando doce meses después.


\subsection{Clasificación económica: Presupuestos Generales del Estado (serie amarilla)}
Aplicada la consolidación intrasectorial que acabamos de fundamentar, con eliminación simétrica de todas las transferencias que cruzan entre institutociones de un mismo sector, el gasto total consolidado de los PGE23 asciende a $553,5$mm de euros ($39,8\%$ del PIB) y los ingresos totales consolidados, sin fondos Next Generation EU, ascienden a $371,8$mm ($26,8\%$ del PIB).

Veamos las composiciones de cada una de las perspectivas.

\subsubsection{Gasto}
El capítulo con más dotación sigue siendo el de transferencias corrientes, con el $61,5\%$ del total. Además, si se añaden las transferencias de capital, el importe conjunto asciende al $64,3\%$. De este volumen total, el $64,5\%$ se destinan a familias e instituciones de lucro, siendo su componente más importante las pensiones. Las transferencias a las Administraciones Territoriales representan el $21,6\%$ del total, la manifestación presupuestaria del sistema de financiación autonómica. Para probar el resultado de la consolidación, solo tenemos que recuperar las cifras expuestas en las anteriores secciones. La agregación de las cuatro entidades por separado ($356,1 + 44,7 + 204,2 + 9,5 = 614,5$mm) frente a los $553,5$mm consolidados arroja una diferencia de $61$mm eliminados por duplicidad, que es la magnitud de las transferencias internas que hemos ido documentando.

Por otro lado, las operaciones no financieras (Capítulos I a VII) suponen $425,6$mm ($30,6\%$ del PIB), la cifra que computa a efectos de saldo público. 

Por último, la clasificación orgánica (subsectores) confirma la jerarquía que veníamos anticipando. Estado $46,8\%$, Seguridad Social $41,9\%$, Organismos Autónomos $9,3\%$, resto de entidades $2\%$. El Estado y la Seguridad Social concentran, entre ambos, casi nueve de cada diez euros del gasto consolidado total.

\subsubsection{Ingresos consolidados: la centralidad de las cotizaciones sociales, reconfirmada desde otro ángulo}
Desde la óptica de ingresos, las fuentes de renta más importantes son los impuestos directos y las cotizaciones sociales, que suponen el $72,4\%$ del total. De estos, los mayores ingresos provienen de las cotizaciones sociales, lo cual es debido a la caja única de la Seguridad Social, que hace que todas las cotizaciones sociales sean ingresadas por la Administración Central, a diferencia de los impuestos.

Las cotizaciones sociales ($184,5$mm, el $49,6\%$) superan, con holgura, tanto a los impuestos directos ($84,5$mm, $22,7\%$) como a los indirectos ($57,8$mm, $15,5\%$). Por tanto, no es solo que las cotizaciones pesen mucho dentro de la presión fiscal agregada de todas las AAPP (como veremos en [\authorfont{Ver Tema 4.B.23}]), sino que dentro del propio perímetro consolidado del Estado y sus entidades dependientes, son literalmente el recurso mayoritario, por encima incluso del conjunto de la imposición directa e indirecta sumadas. 

La razón técnica es la caja única del sistema contributivo, que implica que la totalidad de las cotizaciones sociales devengadas en España llegan, íntegras, al perímetro consolidado que estamos analizando, mientras que una parte sustancial de IRPF, IVA e Impuestos Especiales nunca llega a computar aquí, precisamente porque se cede o se transfiere fuera de este perímetro concreto. 

La necesidad de endeudamiento bruto consolidada asciende a $193,6$mm, el $13,9\%$ del PIB. Es decir, la magnitud total de deuda que el conjunto del sector público administrativo estatal necesita emitir, sumando déficit del ejercicio y refinanciación de vencimientos.

\subsection{Clasificación por políticas de gasto: la funcional, por fin, en su versión consolidada y más reveladora}

Llegamos a la aplicación consolidada de la clasificación funcional. 

El gasto consolidado se organiza en $27$ políticas de gasto, agrupadas en cinco áreas. Considerando el gasto de los capítulos I a VIII. 

Los Servicios Públicos Básicos absorben el 6,0 por ciento del gasto, siendo Defensa y Seguridad ciudadana las pol íticas más destacadas, con un 2, 7 y un 2,3 por ciento, respectivamente, del total.

Las Actuaciones de Protección y Promoción Social absorben el 56 por ciento del gasto total. siendo las políticas más destacadas: Pensiones ( 4 1 .8 por ciento), Otras Prestaciones Económicas (4,9 por ciento) y Desempleo (4,7 por ciento).

La Producción de Bienes Públicos de Carácter Preferente tiene asignado el 2,5 por ciento del gasto, destacando la política de Sanidad con el 1 ,8  por ciento.

Las Actuaciones de Carácter Económico suponen el 7,9 por ciento del gasto. Ocupa el primer lugar Investigación. Desarrollo, Innovación y Digital ización ( 1 ,9 por ciento). seguido de Agricultura, Pesca y Alimentación con el I ,8 por ciento: y de Infraestructuras y Ecosistemas Resilientes ( 1 .8 por ciento).

Las A ctuaciones de Carácter General suponen el 27,6 por ciento del gasto. EntTe ellas están: Transferencias a otras Administraciones Públicas ( 1 4,5 por ciento), Servicios de Carácter General (5,7 por ciento) y Deuda Pública ( 6,9 por ciento).

Conviene señalar que el catálogo salta directamente del Área $4$ a la Área $9$, sin que existan Áreas $5$ a $8$. Es un patrón habitual en los sistemas de clasificación de larga vida, donde categorías antiguas se eliminan o se refunden con el paso de las reformas sucesivas, pero sus números de referencia se dejan deliberadamente vacantes en lugar de renumerar todo el catálogo, precisamente para preservar la comparabilidad de las series históricas con ejercicios presupuestarios anteriores a la reforma. La misma lógica de continuidad estadística que justifica que el INE revise cifras de PIB sin alterar la estructura general de sus series.


E l gasto social es un concepto que se destaca mucho a la hora de evaluar y/o clasificar los presupuestos. En la clasificación por políticas que utiliza la in fonnación presupuestaria. es la suma de las Actuaciones de Protección y Promoción Social y de las de Producción de Bienes Públicos de Carácter Preferente. Por lo tanto, de acuerdo con las cifras consignadas anteriormente, suman un 58,5 por ciento del gasto. Restando las prestaciones por desempleo la cifra queda en el 53,8 por ciento.

La suma del gasto social en sentido amplio —Protección y Promoción Social más Producción de Bienes Públicos de Carácter Preferente— alcanza el 58,5% del gasto consolidado (53,8% si excluimos el desempleo, de naturaleza más cíclica que estructural). [Probable] Esta cifra, leída junto con todo lo que ya construimos en el epígrafe 1.3 sobre la Seguridad Social, permite una lectura de conjunto que conviene explicitar como cierre analítico de todo el tema: más de la mitad de cada euro que gasta el sector público administrativo estatal consolidado se destina, en último término, a sostener a la población de mayor edad o en situación de vulnerabilidad —pensiones contributivas y no contributivas, prestaciones por incapacidad, dependencia, Ingreso Mínimo Vital—, con las pensiones, en solitario, absorbiendo más de cuatro de cada diez euros del gasto total consolidado. Es, precisamente, la magnitud que hace tan urgente y tan políticamente sensible toda la reforma de pensiones que ya trabajamos —Fondo de Reserva, Mecanismo de Equidad Intergeneracional—: cualquier desviación, por pequeña que parezca en términos porcentuales, sobre una partida que representa el 41,8% del gasto total tiene, en términos absolutos, un impacto presupuestario de una magnitud que ninguna otra política de gasto del catálogo puede igualar.

\begin{specialblock}{\textbf{Plan de Recuperación, Transformación y Resilencia:} ¿Qué tratamiento recibió?}
    El PRTR y REACT-EU añaden 30mm adicionales al gasto de 2023 (+6,6% sobre los Capítulos I a VIII), con un reparto que resulta, a primera vista, sorprendente si lo comparamos con el del presupuesto ordinario que acabamos de recorrer: 71,7% se dirige a Actuaciones de Carácter Económico (I+D+i y Digitalización 30,3%, Infraestructuras y Ecosistemas Resilientes 15,1%, Industria y Energía 14,9%), frente a solo un 15,7% a Protección y Promoción Social y un 10,1% a Producción de Bienes Públicos de Carácter Preferente —exactamente lo contrario de la jerarquía que domina el presupuesto ordinario, donde la Protección Social concentraba el 56% frente al 7,9% de las Actuaciones Económicas—. [Seguro] Esta inversión no responde a ninguna decisión discrecional española: es una consecuencia directa y obligatoria del propio Reglamento (UE) 2021/241, de 12 de febrero de 2021, por el que se establece el Mecanismo de Recuperación y Resiliencia, que exige a cada Estado miembro destinar, como mínimo, el 37% del gasto de su plan nacional a objetivos climáticos y el 20% a la transformación digital, sometido además al principio de "no causar un perjuicio significativo" (Do No Significant Harm, DNSH) a ningún objetivo medioambiental, y al principio de adicionalidad —el MRR no puede, salvo excepciones tasadas, sustituir gasto público ya recurrente del Estado miembro—. 
    
    [Seguro] El Plan de Recuperación, Transformación y Resiliencia español, aprobado por la Comisión Europea el 16 de junio de 2021 por un importe de 69.500 millones de euros, superó además el umbral mínimo exigido, alcanzando aproximadamente el 40% de contribución climática frente al 37% exigido. [Probable] Este marco regulatorio europeo explica, con precisión que el documento de referencia no ofrece, por qué el gasto financiado con fondos europeos y el gasto financiado con recursos ordinarios del Estado siguen, dentro del mismo ejercicio presupuestario, lógicas de asignación casi opuestas: el presupuesto ordinario responde a compromisos de gasto social acumulados durante décadas —pensiones, sanidad, dependencia—, mientras que el PRTR responde a una arquitectura de condicionalidad diseñada en Bruselas específicamente para acelerar la transición ecológica y digital de toda la Unión, con independencia de cuáles fueran las prioridades de gasto social ya consolidadas en cada Estado miembro receptor.
\end{specialblock}

FEDEA: prioridades absolutas frente a prioridades relativas

Cierro con la distinción que FEDEA propone para leer estas mismas cifras desde una óptica dinámica y no solo estática: las prioridades absolutas —identificadas por el aumento de peso en el PIB de una partida concreta— señalan a las pensiones y a la financiación de las Comunidades Autónomas como las de mayor peso creciente; las prioridades relativas —identificadas por la tasa de crecimiento anual medio a precios constantes, con independencia de su peso absoluto— señalan, en cambio, a la Dependencia y los Servicios Sociales, la Cultura y el Deporte, las subvenciones al transporte, y la Educación, como las partidas que más rápidamente han crecido en términos relativos, aunque partan de una base mucho más reducida que las pensiones.

Esta doble lectura —qué pesa más en términos absolutos frente a qué crece más deprisa en términos relativos— es, en cierto sentido, el cierre metodológico perfecto de todo este tema: nos recuerda que ninguna cifra presupuestaria, por bien clasificada y bien consolidada que esté, se interpreta correctamente sin preguntarse, además, desde qué óptica temporal y desde qué magnitud de referencia se está leyendo.


















- Ley de Contabilidad y Administración Pública de 1850 (ya citada en el Tema 21).
- Naciones Unidas, Departamento de Asuntos Económicos: A Manual for Economic and Functional Classification of Government Transactions, Nueva York, 1958.
- "La clasificación funcional de los gastos del Estado", Documentación Administrativa, Instituto Nacional de Administración Pública (nota informativa sobre la recepción española del Manual de la ONU de 1958).
- CEPAL, ILPES y DOAT-ONU: Esquema Uniforme de Clasificación de las Cuentas del Sector Público adaptado a los países de América Latina, Grupo de Trabajo sobre Contabilidad Fiscal Uniforme, Santiago de Chile, noviembre de 1964.
- Seminarios de Clasificación y Administración Presupuestarias en Sudamérica (Santiago de Chile, 1962) y en Centroamérica y Panamá (San José de Costa Rica, 1963).
- WILDAVSKY, A. (1964): The Politics of the Budgetary Process (ya citado en el Tema 21, aplicado ahora al origen específico de la presupuestación por programas española).
- Doctrina de la responsabilidad ministerial individual (individual ministerial responsibility), tradición constitucional británica de Westminster.
- STONE, R. y MEADE, J. (1941): An Analysis of the Sources of War Finance and an Estimate of the National Income and Expenditure in 1938 and 1940, White Paper del Gobierno británico.
- STONE, R. (1953 y 1968): A System of National Accounts, ONU (bases del Sistema de Cuentas Nacionales).
- Real Academia Sueca de Ciencias: comunicado del Premio Nobel de Economía 1984, "por sus contribuciones fundamentales al desarrollo de los sistemas de cuentas nacionales".
- DEATON, A. (2008): "Stone, John Richard Nicholas (1913-1991)", The New Palgrave Dictionary of Economics.
- Naciones Unidas (1958): A Manual for Economic and Functional Classification of Government Transactions (ya citado en el epígrafe 1.1).
- Literatura de Nueva Gestión Pública sobre presupuestación por resultados (outputs/outcomes) frente a presupuestación departamental clásica: casos de Nueva Zelanda, Australia y Reino Unido.
- Real Decreto Legislativo 2/2004, de 5 de marzo, Texto Refundido de la Ley Reguladora de las Haciendas Locales, art. 167.1 y disposición adicional sexta.
- Orden del Ministerio de Economía y Hacienda de 20 de septiembre de 1989 (primera estructura presupuestaria local obligatoria, vigente desde 1992).
- Orden EHA/1645/2004, de 3 de junio, por la que se dictan las normas para la elaboración de los Presupuestos Generales del Estado para 2005 (introducción de la clasificación por programas en sustitución de la funcional).
- Orden EHA/3565/2008, de 3 de diciembre, por la que se aprueba la estructura de los presupuestos de las entidades locales (BOE núm. 297, de 10 de diciembre de 2008).
- Orden HAP/419/2014, de 14 de marzo, por la que se modifica la Orden EHA/3565/2008.
- Leyes de Hacienda Pública de las distintas Comunidades Autónomas de régimen común.
- Ley 12/2002, de 23 de mayo (Concierto Económico con el País Vasco) y Ley 28/1990, de 26 de diciembre (Convenio Económico con Navarra), ya citadas en el Tema 21.
- Ministerio de Hacienda: Central de Información Económico-Financiera de las Administraciones Públicas.
- PGE 2023, Informe Económico y Financiero (fuente principal de cifras, ya citada).
- Constitución Española, art. 135.3 (prioridad absoluta de pago de la deuda pública), ya trabajado en el Tema 21.
- Ley Orgánica 2/2012, art. 30.1 (fórmula de cálculo del techo de gasto no financiero).
- Boletín Oficial de las Cortes Generales (2020): Informe de evaluación y reforma del Pacto de Toledo
- Congressional Budget Office / Office of Management and Budget (EEUU): distinción entre mandatory y discretionary spending, como referencia comparada del fenómeno de estrechamiento del margen de gasto discrecional.
- Tesoro Público: Estrategia de emisión de Deuda del Estado, ejercicio 2023 (emisión neta y bruta).
- Real Decreto-ley 36/1978, de 16 de noviembre, de creación del Instituto Nacional de Empleo.
- Ley 56/2003, de 16 de diciembre, de Empleo, disposición adicional primera (cambio de denominación de INEM a SEPE y creación del Sistema Nacional de Empleo).
- Ley 3/2023, de 28 de febrero, de Empleo (propuesta de transformación del SEPE en Agencia Estatal).
- RODRÍGUEZ-PIÑERO, M. (1994): sobre el fin del monopolio público de colocación del INEM.
- Reglamento (UE) núm. 1307/2013 y sucesivos, sobre pagos directos de la PAC (FEAGA); Reglamento (UE) núm. 1305/2013 y sucesivos, sobre desarrollo rural (FEADER).
- Reforma MacSharry de la Política Agrícola Común, 1992.
- Directiva 2008/94/CE del Parlamento Europeo y del Consejo, de 22 de octubre de 2008, sobre la protección de los trabajadores asalariados en caso de insolvencia del empresario (base del FOGASA).
- Real Decreto 375/2003, de 28 de marzo, por el que se aprueba el Reglamento General de Mutualismo Administrativo (MUFACE).
- Instituto Nacional de Previsión, creado en 1908; Ley de Bases de la Seguridad Social de 1963 y Texto Articulado de 1966 (Decreto 907/1966).
- Constitución Española de 1978, art. 41.
- Decreto Ley 3.500, de 1980 (Chile), creación del sistema de Administradoras de Fondos de Pensiones (AFP) bajo impulso de José Piñera.
- Ley 24/1997, de 15 de julio, de Consolidación y Racionalización del Sistema de Seguridad Social (creación del Fondo de Reserva).
- Ley 28/2003, de 29 de septiembre, reguladora del Fondo de Reserva de la Seguridad Social.
- Pacto de Toledo, 1995 (origen del consenso sobre la constitución de un fondo de reserva).
- Ley 21/2021, de 28 de diciembre, de garantía del poder adquisitivo de las pensiones y de otras medidas de refuerzo de la sostenibilidad financiera y social del sistema público de pensiones.
- Real Decreto-ley 2/2023, de 16 de marzo, de medidas urgentes para la ampliación de derechos de los pensionistas, la reducción de la brecha de género y el establecimiento de un nuevo sistema de cotización para los trabajadores por cuenta propia.
- Real Decreto-ley 13/2022, de 26 de junio (reforma del sistema de cotización de trabajadores autónomos por tramos de rendimientos netos).
- Ley 3/2013, de 4 de junio, de creación de la Comisión Nacional de los Mercados y la Competencia.
- Real Decreto 657/2013, de 30 de agosto, por el que se aprueba el Estatuto Orgánico de la CNMC.
- Auto del Tribunal Supremo (Sala de lo Contencioso-Administrativo, Sección 3ª) 5896/2015, de 3 de julio (cuestiones prejudiciales ante el TJUE sobre la CNMC).
- Directiva 2002/21/CE y Directiva 2009/140/CE (marco regulador común de las redes y los servicios de comunicaciones electrónicas).
- Real Decreto 1067/2015, de 27 de noviembre, por el que se aprueba el Estatuto de la Agencia Estatal de Investigación.
- Ley 31/1990, de 27 de diciembre, de Presupuestos Generales del Estado para 1991 (creación de la AEAT).
- Real Decreto de 24 de enero de 1907, de creación de la Junta para Ampliación de Estudios e Investigaciones Científicas.
- Ley de 24 de noviembre de 1939, de creación del Consejo Superior de Investigaciones Científicas.
- Ley 23/1998, de 7 de julio, de Cooperación Internacional para el Desarrollo.
- Naciones Unidas: objetivo del 0,7% de la Renta Nacional Bruta en ayuda oficial al desarrollo (Resolución 2626 (XXV) de la Asamblea General, 1970).
- Fondo Monetario Internacional: Government Finance Statistics Manual (GFSM), ediciones 1986, 2001 y 2014, en especial los párrafos 3.155 a 3.157 de la edición de 2014 (consolidación intrasectorial e intersectorial).
- O'CONNOR, K. W., WEISMAN, E. y WICKENS, T. (2004): "Consolidation of the General Government Sector", Government Finance Statistics Manual 2001, Companion Material, Fondo Monetario Internacional.
- Oficina de Estadística de Australia (Australian Bureau of Statistics): Australian System of Government Finance Statistics: Concepts, Sources and Methods, Parte H (Consolidación), edición 2015.

- PGE 2023, Informe Económico y Financiero, cifras consolidadas (ya citado).
- Reglamento (UE) 2021/241 del Parlamento Europeo y del Consejo, de 12 de febrero de 2021, por el que se establece el Mecanismo de Recuperación y Resiliencia, en especial arts. 5 (adicionalidad), 18 y Anexo VI (umbrales mínimos de gasto climático y digital), y art. 5.2 (principio de no causar un perjuicio significativo, DNSH).
- Decisión de Ejecución del Consejo relativa a la aprobación del Plan de Recuperación, Transformación y Resiliencia de España, 16 de junio de 2021.
- Orden HFP/1030/2021 y Orden HFP/1031/2021, de 29 de septiembre (sistema de gestión del PRTR).
- FEDEA (2022.a): Notas sobre los Presupuestos Generales del Estado de 2023, Apuntes 22/24 (ya citado en el Tema 22, prioridades absolutas y relativas del gasto).












































\clearpage
\section{Conclusión} 

La Administración Central moviliza anualmente unos recursos equivalentes al 40 por ciento del PIB. según su presupuesto consolidado. El gasto financiado con fondos Next Generation EU incrementan esta cifra algo más de un 5 por ciento. A estos recursos hay que aiiadir los de sociedades y entidades con presupuesto estimativo. los cuales son de menor importancia cuantitativa, pero esenciales para la producción de bienes y serv icios. para la supervisión del sistema financiero y para la provisión de crédito y seguro de crédito a la economía. Es el caso. por ��iemplo, de cie11as sociedades y entidades mercantiles. que juegan un papel crucial en la actividad inversora del sector público estat¡.¡I que, en total, asciende casi al 2 por ciento del Pll3.

En el entramado de los PGE juega un papel central el Estado, con el 46,8 por ciento del presupuesto, seguido de la Seguridad Social, con el 41.9 por ciento. Aparte de la cuantía, el carácter cemral del presupuesto del Estado se refleja en que, por medio de transferencias corrientes y de capital. suministra recursos al resto de administraciones públicas. sin los cuales éstas no podrían operar. También proporciona recursos al resto de la econom ía, pues a lo largo del tiempo, por medio de transferencias. el Estado financia políticas de todo tipo: subvenciones al coste de la energía. becas. subvenciones al transporte. compensaciones a concesionarias de autopistas de peaje·. apoyo a la atención san itaria etc.

A pesar del gran volumen de recursos movilizados, hay que destacar el escaso margen para emprender políticas nuevas. Efectivamente. gran pai1e del gasto está vinculado a la legislación. Otra gran parte corresponde a políticas que llevan décadas implementándose. las cuales se pueden redisef'iar. pero difícilmente elim inar.

Dado que en estos presupuestos se incluyen organismos como la Seguridad Social y el SEPE, va de suyo que el gasto social es el que más recursos absorbe. con el 58.5 por ciento (el 53.8 por ciento sin contar el gasto de desempleo); que es la suma de las Actuaciones de Protección y Promoción Social y la Producción de Bienes Públicos de Carácter Preferente. Es una cifra bastante considerable teniendo en cuenta que el Estado apenas tiene competencias en sanidad o educación.

En cambio, proporciona financiación a las administraciones territoriales para estas y otrns políticas.

Por el lado de los ingresos, la administración central obtiene el 72 por ciento de sus recursos de impuestos directos y cotizaciones sociales y un 16 por ciento de impuestos indirectos. Además, necesita recurrir a un endeudamiento bruto equivalente al 52, 1 por ciento de la suma de ingresos de capítulo la capítulo VII 1, sin fondos Next Generation EU. Estos fondos aportan un 5 por ciento adicional a la recaudación. algo similar a lo que ocurre con los gastos.40


\subsection{Recapitulación}


\subsection{Extensiones y relación con otras partes del Temario}



\subsection{Opinión}

\subsubsection{Idea final (Salida o cierra)}

\clearpage
\pagenumbering{gobble}
\MyBibliography



\MyBibItem{DAWID y GATTI}{Chapter 2 - Agent-Based Macroeconomics}{2018}{https://doi.org/10.1016/bs.hescom.2018.02.006}


% ANEXOS
\clearpage
\anexo{Material de complemento: Entidades de presupuesto estimativo}\label{anx:presupuesto_estimativo}
He investigado con el mismo rigor que el resto del tema, incluyendo el episodio más actual y más delicado de todos —el FASEE y el caso Plus Ultra—, que trato con la neutralidad que exige tratarse de un procedimiento judicial en curso: expongo los hechos verificables y las distintas posiciones, sin pronunciarme sobre responsabilidades que corresponde determinar a los tribunales.

---

## Anexo. Las entidades con presupuesto estimativo

### Fundamento teórico: por qué existen, y por qué el enfoque clasificatorio de este tema no puede aplicárseles

[Seguro] Ya fundamentamos, en el Bloque 1, la razón última de esta frontera: la clasificación orgánica/económica/funcional es el lenguaje técnico que hace posible el principio de especificación —el límite legal cuantitativo y cualitativo que el Parlamento impone al gasto—, y ese lenguaje solo tiene sentido cuando existe, efectivamente, un límite que vigilar. [Probable] Las entidades que reunimos en este anexo —sociedades mercantiles estatales, entidades públicas empresariales— desarrollan, por definición, una actividad de naturaleza económica, con ingresos propios derivados de su actividad en el mercado (venta de billetes, cánones de navegación aérea, tasas portuarias, intereses de crédito), y su presupuesto no es, por tanto, un techo legal que las Cortes fijan y cuyo respeto haya que verificar partida por partida, sino una previsión de explotación —análoga a la cuenta de pérdidas y ganancias de cualquier empresa— y una previsión de capital —análoga a un estado de flujos de efectivo—, ambas de naturaleza estimativa: pueden superarse sin que ello constituya, por sí solo, ninguna irregularidad, precisamente porque no protegen ningún límite parlamentario específico sino que anticipan, con la incertidumbre propia de cualquier previsión empresarial, cómo evolucionará su actividad económica a lo largo del ejercicio. [Suposición] Esta imposibilidad de encaje no es, sin embargo, una simple curiosidad técnica: es, como veremos, la razón por la que estas entidades han sido, históricamente, el terreno más fértil para los grandes debates de política económica española —privatización, liberalización de mercados, rescates empresariales— que la parte "administrativa" del sector público, mucho más rígida y mucho menos discrecional, raramente protagoniza con la misma intensidad.

### El fundamento político-económico de fondo: la remisión necesaria al Tema 15, y lo que este anexo añade

[Probable] El debate de fondo sobre por qué debe existir, o no, la empresa pública —Vernon, Baumol-Panzar-Willig, la escuela de la elección pública de Buchanan y Tullock, Niskanen, Kornai y la restricción presupuestaria blanda, Hart-Shleifer-Vishny, Mazzucato y el Estado emprendedor— ya lo construimos con la profundidad que merece en el Tema 4.B.15, y no tiene sentido repetirlo aquí sin más. Lo que este anexo aporta, que aquel tema no podía ofrecer, es la aplicación presupuestaria concreta de ese debate a las entidades reales que hoy componen el sector público empresarial estatal español, con sus cifras, sus controversias más recientes y su situación europea actualizada.

### SEPI: de la herencia del INI a la gestión del fondo de rescate más controvertido de la última década

[Probable] La Sociedad Estatal de Participaciones Industriales, creada en 1995, es la heredera directa del Instituto Nacional de Industria (INI), fundado en 1941 bajo el franquismo como instrumento central de la política de industrialización autárquica del régimen, y del Instituto Nacional de Hidrocarburos —la misma genealogía de empresa pública española que ya trabajamos con todo detalle en el Tema 15, ahora con su desenlace institucional concreto—. Su cartera actual incluye participaciones en Navantia, Tragsa, Correos, Hulleras del Norte y, temporalmente, en las empresas rescatadas durante la pandemia. [Seguro] Y es precisamente en este último papel donde SEPI ha protagonizado el episodio más controvertido de toda la actualidad reciente del sector público empresarial español: el Fondo de Apoyo a la Solvencia de Empresas Estratégicas (FASEE), creado en julio de 2020 y dotado con 10.000 millones de euros, gestionado por SEPI para apoyar, mediante préstamos participativos y ordinarios —no mediante ayudas a fondo perdido—, a empresas consideradas estratégicas afectadas por la pandemia. [Seguro] De las 74 solicitudes recibidas, por un importe conjunto de 5.393 millones, se aprobaron finalmente 30 operaciones por 3.256 millones; entre ellas, el préstamo de 475 millones a Air Europa —el mayor del fondo— y el de 53 millones a la aerolínea Plus Ultra, esta última sometida, desde entonces, a un intenso escrutinio público y judicial. [Seguro] La operación Plus Ultra generó, desde el primer momento, dudas parlamentarias sobre su carácter "estratégico" —una aerolínea con una plantilla inferior a 300 personas y con parte de su capital vinculado a inversores venezolanos—, dudas que derivaron, con el tiempo, en una investigación judicial en la Audiencia Nacional que afecta al expresidente del Gobierno José Luis Rodríguez Zapatero, por presuntos delitos de blanqueo de capitales y tráfico de influencias en relación con su papel de intermediación en la concesión de la ayuda; el propio rescate de Air Europa ha quedado, además, vinculado en la instrucción judicial al llamado "caso Koldo". [Seguro] Conviene, con el mismo rigor que hemos aplicado a cada controversia de este proyecto, presentar también el otro lado de este episodio: el Tribunal de Cuentas, en su fiscalización de la actuación de SEPI en el procedimiento FASEE, concluyó que la entidad "aplicó adecuadamente las disposiciones legales en materia de cumplimiento de los requisitos de elegibilidad" y que "las condiciones de la financiación otorgada... se adecuaron a la legalidad", sin apreciar perjuicio contable alguno según el propio informe de la Intervención General del Estado. [Probable] A la fecha más reciente disponible, el fondo ha recuperado ya el 65% de los préstamos concedidos —1.745 de los 2.681 millones financiados a las 28 empresas rescatadas, con once de ellas habiendo devuelto ya la totalidad de su ayuda—, aunque diecisiete empresas, entre ellas la propia Plus Ultra y Tubos Reunidos, siguen atravesando dificultades que han obligado a SEPI a renegociar sus calendarios de pago. Es, en su conjunto, un episodio que ilustra con extraordinaria nitidez el riesgo que la teoría del riesgo moral y de la restricción presupuestaria blanda de Kornai —ya trabajada en el Tema 15— anticipa siempre que el Estado se convierte en prestamista de última instancia de empresas privadas en dificultad: la propia existencia del fondo, necesaria para sostener el tejido productivo durante una crisis excepcional, genera simultáneamente el riesgo de que su gestión discrecional —qué empresa es "estratégica" y cuál no— se convierta en terreno fértil para la influencia política inadecuada, exactamente el temor que la literatura académica sobre empresa pública lleva décadas señalando en abstracto y que este episodio concreto ha convertido, para España, en un caso de estudio con nombre y apellido propio.

### ENAIRE y AENA: la privatización parcial que convive con la propiedad pública mayoritaria

[Seguro] ENAIRE, gestor de la navegación aérea, mantiene naturaleza íntegramente pública, pero conviene distinguirlo con precisión de AENA, el gestor de los aeropuertos españoles, porque su situación jurídica es distinta y revela un modelo intermedio entre la propiedad pública plena y la privatización completa que rara vez se explica con claridad. El Gobierno decidió, en julio de 2014, la privatización parcial de AENA, materializada mediante su salida a Bolsa el 11 de febrero de 2015 —trece meses después de la decisión—, colocando en el mercado una participación minoritaria mientras el Estado, a través de ENAIRE, conservaba el 51% del capital y, con él, el control efectivo de la compañía. [Probable] Es un modelo de privatización parcial y deliberadamente limitada —muy distinto de las privatizaciones íntegras que ya trabajamos en el Tema 15 con el caso británico thatcherista—, que permite al Estado captar capital privado y disciplina de mercado sin renunciar al control estratégico de una infraestructura considerada esencial; el propio folleto de la Oferta Pública de Venta de AENA reconocía, de forma reveladora, a la competencia del AVE como uno de los riesgos de negocio más relevantes para sus futuros accionistas —un dato que conecta, de forma inesperada, con el siguiente apartado de este mismo anexo.

### ADIF, ADIF-Alta Velocidad y RENFE: de un monopolio centenario a uno de los mercados ferroviarios más abiertos de Europa

[Seguro] El sector ferroviario español ha protagonizado, en los últimos cinco años, una de las transformaciones más rápidas y más vigiladas internacionalmente de todo el sector público empresarial español. La separación entre ADIF —gestor de la infraestructura ferroviaria— y RENFE-Operadora —operador de trenes— responde, en origen, a una exigencia de las Directivas europeas de liberalización ferroviaria, que exigen separar la gestión de la infraestructura de la prestación del servicio precisamente para permitir la entrada de operadores competidores sin que el antiguo monopolista controle, al mismo tiempo, las propias vías por las que sus competidores deben circular —el mismo principio de separación vertical que ya mencionamos, en otro contexto, al hablar del origen de la CNMC en el Bloque 2—. [Seguro] Esa liberalización, ordenada desde Bruselas, se materializó en España con una rapidez que ha convertido al país en una referencia europea: OUIGO, filial de la francesa SNCF, inició sus servicios entre Madrid y Barcelona el 11 de mayo de 2021; IRYO, consorcio de Air Nostrum y la italiana Trenitalia, entró en el mismo corredor el 25 de noviembre de 2022. [Seguro] El resultado, documentado por la propia CNMC, es que España se convirtió en el primer país de Europa en lograr la competencia efectiva de tres operadores en un mismo corredor de alta velocidad —Renfe (con sus dos marcas, AVE y el low-cost AVLO), Ouigo e Iryo—, con un impacto medible sobre el mercado: los viajeros de alta velocidad crecieron un 37% en 2023 respecto al año anterior, hasta 32,4 millones, mientras la cuota de mercado de Renfe caía hasta diez puntos porcentuales en los corredores con competencia —situándose en torno al 58% en Madrid-Barcelona, frente al 22% de Iryo y el 20% de Ouigo—. [Probable] El episodio tiene, además, un epílogo de rabiosa actualidad que conviene señalar como cierre irónico: mientras España abrió sus corredores más rentables a la competencia europea con una celeridad que sus propios competidores extranjeros reconocen como excepcional, Francia —el país de origen de Ouigo— ha puesto, según reportan varios medios especializados, notables dificultades administrativas a que la propia Renfe pueda, en reciprocidad, operar en territorio francés, hasta el punto de que la operadora española estaría valorando abandonar sus intentos de entrada en el mercado galo — un recordatorio de que la liberalización europea del transporte ferroviario, pese a su marco jurídico común, no se aplica siempre con la misma reciprocidad efectiva en la práctica de cada Estado miembro.

### El ICO y el resto del catálogo: brevedad justificada por remisión a lo ya trabajado

[Probable] El Instituto de Crédito Oficial, banco público de desarrollo español, desempeñó durante la pandemia un papel de primer orden mediante sus líneas de avales a la financiación empresarial —el mismo mecanismo de pasivo contingente que ya trabajamos con detalle en el Tema 21, a propósito de los 14,5 puntos de PIB en avales concedidos por el Estado a finales de 2022—, sin que sea necesario redesarrollar aquí esa misma pieza. CESCE, el Consorcio de Compensación de Seguros, las fundaciones del sector público, y el propio Banco de España —con su propuesta de gastos remitida anualmente a las Cortes conforme a la Ley de Autonomía de 1994— completan un catálogo ya suficientemente trabajado, en sus cifras concretas, dentro del cuerpo principal del Tema 22.

### La referencia internacional que ordena todo este catálogo: las directrices de la OCDE sobre gobierno corporativo de las empresas públicas

[Seguro] Cierro con el marco de referencia internacional que da coherencia comparada a todo este anexo: las Directrices de la OCDE sobre el Gobierno Corporativo de las Empresas Públicas, actualizadas por última vez en 2015, constituyen el estándar internacional de referencia para cómo un Estado debe organizar, gobernar y rendir cuentas de su cartera de empresas públicas —recomendando, entre otros principios, una separación clara entre la función del Estado como propietario y su función como regulador del mercado en que esa empresa opera (precisamente la separación que justificó, en su día, la creación de la CNMC como regulador independiente distinto del propio accionista estatal), y una publicación regular y agregada de la cartera empresarial pública comparable a la de cualquier inversor institucional privado—. [Probable] Y sobre las operaciones más delicadas de todo este anexo —las inyecciones de capital público a empresas en dificultades, como las del propio FASEE—, el Derecho de la competencia de la Unión Europea somete a cualquier Estado miembro al llamado criterio del inversor privado (private investor test): la Comisión Europea evalúa si un inversor privado, en las mismas circunstancias de mercado, habría realizado la misma operación en los mismos términos: solo si la respuesta es afirmativa la operación escapa a la calificación de ayuda de Estado prohibida por el art. 107 del TFUE, sometida en caso contrario a la obligación de notificación previa y autorización de la propia Comisión. Es, en definitiva, el mismo tipo de escrutinio europeo —previo, externo, técnico— que ya hemos visto operar, con otros objetos y en otros temas de esta serie, sobre prácticamente cualquier decisión de calado del sector público español: ni siquiera el terreno, aparentemente más discrecional y más alejado de la disciplina presupuestaria clásica, de las entidades con presupuesto estimativo, escapa por completo a la vigilancia de Bruselas.

### Verificación final

Repaso lo cubierto: el fundamento teórico de por qué el aparato clasificatorio de este tema no puede aplicarse a estas entidades, con remisión explícita y no redundante al aparato teórico completo ya construido en el Tema 15; SEPI, con su genealogía desde el INI franquista hasta el episodio FASEE-Plus Ultra tratado con la neutralidad que exige un procedimiento judicial en curso; ENAIRE y AENA como ejemplo de privatización parcial y no completa; ADIF y RENFE con la liberalización ferroviaria europea y sus cifras de competencia más actualizadas, incluido el epílogo de reciprocidad franco-española; el resto del catálogo, con la brevedad justificada por su ya amplio tratamiento previo; y el doble marco de referencia internacional que ordena todo el conjunto —las directrices de gobierno corporativo de la OCDE, y el criterio del inversor privado del Derecho europeo de la competencia—. Con esto, y dado que este anexo es, como bien señalaste, el único lugar de todo el temario donde estas entidades reciben tratamiento propio, considero completo el desarrollo que la finalidad del Tema 22 exigía.



- Ley 5/1996, de 10 de enero, de creación de determinadas entidades de derecho público (creación de SEPI).
- Real Decreto-ley 25/2020, de 3 de julio, de medidas urgentes para apoyar la reactivación económica y el empleo (creación del FASEE).
- Orden PCM/679/2020, de 23 de julio (desarrollo del funcionamiento del FASEE).
- Tribunal de Cuentas: informe de fiscalización de la actuación de SEPI en el procedimiento de aprobación de operaciones del FASEE (remitido al Congreso de los Diputados).
- Comisión Nacional de los Mercados y la Competencia: Informe Anual del Sector Ferroviario, ediciones sucesivas.
- Acuerdos marco de reserva de capacidad entre ADIF-Alta Velocidad y Renfe Viajeros, Iryo y Ouigo, firmados el 13 de mayo de 2020.
- Directiva 2012/34/UE del Parlamento Europeo y del Consejo, de 21 de noviembre de 2012, por la que se establece un espacio ferroviario europeo único (base de la liberalización).
- OCDE (2015): Directrices de la OCDE sobre el Gobierno Corporativo de las Empresas Públicas.
- Art. 107 del Tratado de Funcionamiento de la Unión Europea (ayudas de Estado) y doctrina del criterio del inversor privado en una economía de mercado (Market Economy Investor Principle).

Con esto queda completo el Tema 4.B.22 en su totalidad, incluido este anexo. ¿Quieres que revisemos el tema completo de conjunto, o empezamos ya con la estructura del Tema 4.B.23?









\end{document}